\documentclass[a4paper,10pt]{article}
\usepackage{amsmath,amssymb}
\usepackage{fontspec}
\setmainfont{Liberation Serif}
\usepackage{graphicx}
\usepackage{eso-pic}
\usepackage{tikz}
\usepackage[most]{tcolorbox}
\usepackage{geometry}
\usepackage{fancyhdr}
\usepackage{lastpage}
\usepackage{booktabs}
\usepackage{listings}
\usepackage{enumitem}
\usepackage{microtype}

\geometry{
    a4paper,
    top=1.5cm,
    bottom=1.8cm,
    left=1.2cm,
    right=1.2cm
}

% Persistent Unremovable Canvas Watermark on Every Page
\newsavebox{\watermarkbox}
\sbox{\watermarkbox}{%
    \begin{tabular}{c}
        \includegraphics[width=8cm]{/home/johan/GATE/GATE_CS_PYQ_DB/assets/studiva-butterfly-icon.jpg} \\[-0.5em]
        {\fontsize{18}{22}\selectfont\bfseries\scshape www.studiva.co.in}
    \end{tabular}%
}
\AddToShipoutPictureBG{%
    \begin{tikzpicture}[remember picture, overlay]
        \node[opacity=0.12] at (current page.center) {\usebox{\watermarkbox}};
    \end{tikzpicture}%
}

% Header and Footer
\pagestyle{fancy}
\fancyhf{}
\fancyhead[L]{\small\scshape GATE 2018--2025 CS/IT $\cdot$ \textbf{Programming and Data Structures}}
\fancyhead[R]{\small\scshape Practice Question Bank with Solutions}
\fancyfoot[L]{\scriptsize\textsc{Copyright \copyright\ 2026} Studiva\textsuperscript{TM}}
\fancyfoot[C]{\small Page \thepage\ of \pageref{LastPage}}
\fancyfoot[R]{\scriptsize\texttt{www.studiva.co.in}}
\renewcommand{\headrulewidth}{0.4pt}
\renewcommand{\footrulewidth}{0.4pt}

% Code Listing Style (Monochrome)
\lstset{
    basicstyle=\footnotesize\ttfamily,
    numbers=none,
    frame=none,
    tabsize=2,
    breaklines=true,
    showstringspaces=false
}

\begin{document}

% ==================== PAGE 1: IMPRINT & TITLE PAGE ====================
\thispagestyle{empty}
\vspace*{0.5cm}
\begin{center}
    \includegraphics[width=2.5cm]{/home/johan/GATE/GATE_CS_PYQ_DB/assets/studiva-butterfly-icon.jpg}\\[0.6em]
    {\LARGE \textbf{GATE 2018--2025 COMPUTER SCIENCE \& IT}}\\[0.3em]
    {\Large \textbf{PREVIOUS YEAR QUESTION BANK WITH SOLUTIONS}}\\[0.4em]
    {\large \textsc{\textbf{Subject Wise: Programming and Data Structures}}}\\[0.2em]
    {\normalsize \textsc{Combined 2018--2025 Examination Papers $\cdot$ Comprehensive Edition ($\mathbf{43}$ Questions)}}\\[1.2em]
    \rule{0.6\textwidth}{0.6pt}
\end{center}

\vspace{1em}
\noindent\begin{tabular}{ll}
\textbf{Author / Compiler:} & Rupam Ghosh \\
\textbf{Publisher \& Platform:} & Studiva\textsuperscript{TM} (\texttt{www.studiva.co.in}) \\
\textbf{Official Inquiries:} & \texttt{contact@studiva.co.in} \\
\textbf{Edition:} & First Standard Edition (2026) \\
\textbf{Coverage:} & Complete GATE CS/IT Papers (2018, 2022, 2023, 2024 CS1, 2024 CS2, 2025 CS1, 2025 CS2) \\
\textbf{Syllabus Reference:} & Official GATE 2026 Syllabus (Organizing Institute: IIT Guwahati) \\
\end{tabular}

\vspace{1.5em}
\begin{tcolorbox}[enhanced, standard jigsaw, opacityback=0, colframe=black!40, arc=1mm, boxrule=0.5pt, left=10pt, right=10pt, top=8pt, bottom=8pt]
\textbf{\large COPYRIGHT \& INTELLECTUAL PROPERTY NOTICE}\\[0.6em]
Copyright \copyright\ 2026 Rupam Ghosh \& Studiva\textsuperscript{TM}. All Rights Reserved.\\[0.5em]
This educational compilation, problem organization, taxonomy classification, LaTeX typesetting, hints, and explanatory solutions are the proprietary intellectual property of \textbf{Rupam Ghosh} and \textbf{Studiva\textsuperscript{TM}}.\\[0.6em]
\textbf{Terms of Use \& Attribution Rules:}
\begin{itemize}[leftmargin=*, nosep]
    \item \textbf{Personal Educational Use:} Permitted freely for individual students and educators preparing for competitive examinations.
    \item \textbf{Mandatory Attribution:} Any reproduction, republication, hosting, or sharing (digital or physical) must prominently credit \textbf{Rupam Ghosh \& Studiva\textsuperscript{TM}} and include an active, direct link to \texttt{https://www.studiva.co.in}.
    \item \textbf{Commercial Restrictions:} Commercial resale, sublicensing, course bundling, or inclusion in paid test series without prior written permission is strictly prohibited.
    \item \textbf{Anti-Tampering Notice:} Removing, altering, or obscuring copyright notices, page watermarks, or website links from this document is strictly prohibited.
\end{itemize}
\vspace{0.6em}
\textbf{Fair Use Disclaimer:}\\
The underlying examination questions originate from the Graduate Aptitude Test in Engineering (GATE), conducted by the Indian Institutes of Technology (IITs) and Indian Institute of Science (IISc). Question texts are reproduced for non-commercial educational study under Fair Use doctrines. No copyright is claimed over the original GATE examination question texts.
\end{tcolorbox}

\vspace{1.2em}
\begin{center}
\small\itshape
Practice Recommendation: Use an opaque sheet of paper or ruler to cover the right-hand column during mock attempts. Tick your chosen answers directly in the check boxes on the left, then slide the cover to check the solution and hint.
\end{center}

\newpage
% ==================== PAGE 2 ONWARDS: QUESTION BANK ====================

\noindent
\begin{tcolorbox}[
    blanker,
    sidebyside,
    sidebyside align=top,
    lefthand width=0.55\textwidth,
    top=4pt, bottom=4pt, left=0pt, right=0pt,
    middle=16pt,
    underlay={
        \draw[dashed, black!80, line width=0.8pt, dash pattern=on 4pt off 3pt]
            ([xshift=0.55\textwidth+8pt]frame.north west) --
            ([xshift=0.55\textwidth+8pt]frame.south west);
    }
]
\textbf{Q.1} \quad {\footnotesize\textsc{[MCQ $\cdot$ 1 Mark(s) $\cdot$ Neg: 0.33]}} \hfill {\footnotesize\textbf{GATE 2018 CS} \quad \tiny\texttt{2018\_CS\_GA\_Q02\_v2}}
\vspace{2pt}\hrule height 0.3pt\vspace{4pt}

Consider the following C program.
\begin{lstlisting}[language=C]
#include <stdio.h>
struct Ournode {
    char x, y, z;
};
\par\medskip
int main() {
    struct Ournode p = {'1', '0', 'a' + 2};
    struct Ournode *q = &p;
    printf("%c, %c", *((char*)q + 1), *((char*)q + 2));
    return 0;
}
\end{lstlisting}
The output of this program is

\par\vspace{6pt}\noindent
\begin{tabular}{@{}p{0.48\linewidth}@{\hspace{0.04\linewidth}}p{0.48\linewidth}@{}}
$\square$\ \textbf{(A)} 0, c & $\square$\ \textbf{(B)} 0, a+2 \\
\addlinespace[3pt]
$\square$\ \textbf{(C)} '0', 'p+2' & $\square$\ \textbf{(D)} 1, 0
\end{tabular}

\tcblower

\textbf{Official Key:} \textbf{\large A} \hfill \textit{\small [Easy]}
\vspace{2pt}\hrule height 0.3pt\vspace{3pt}
\textbf{Topic:} Programming and Data Structures \\
\textbf{Subtopic:} \textit{Programming in C}
\vspace{2pt}\hrule height 0.3pt\vspace{3pt}
\textbf{Guide / Hint:} \\
{\small In `struct Ournode`, `char x, y, z;` are stored contiguously in memory. Since `p = {'1', '0', 'a'+2}`, `p.x = '1'`, `p.y = '0'`, and `p.z = 'c'`. Pointer `q` points to `p`; casting `q` to `(char*)` gives byte-level access where `*((char*)q + 1)` accesses `p.y` ('0') and `*((char*)q + 2)` accesses `p.z` ('c').}

\end{tcolorbox}
\vspace{6pt}\hrule height 0.35pt\vspace{6pt}
\noindent
\begin{tcolorbox}[
    blanker,
    sidebyside,
    sidebyside align=top,
    lefthand width=0.55\textwidth,
    top=4pt, bottom=4pt, left=0pt, right=0pt,
    middle=16pt,
    underlay={
        \draw[dashed, black!80, line width=0.8pt, dash pattern=on 4pt off 3pt]
            ([xshift=0.55\textwidth+8pt]frame.north west) --
            ([xshift=0.55\textwidth+8pt]frame.south west);
    }
]
\textbf{Q.2} \quad {\footnotesize\textsc{[NAT $\cdot$ 1 Mark(s) $\cdot$ No Neg]}} \hfill {\footnotesize\textbf{GATE 2018 CS} \quad \tiny\texttt{2018\_CS\_GA\_Q20}}
\vspace{2pt}\hrule height 0.3pt\vspace{4pt}

The postorder traversal of a binary tree is $8, 9, 6, 7, 4, 5, 2, 3, 1$. The inorder traversal of the same tree is $8, 6, 9, 4, 7, 2, 5, 1, 3$. The height of a tree is the length of the longest path from the root to any leaf (in terms of number of edges). The height of the binary tree is \_\_\_\_\_\_\_\_.

\par\vspace{6pt}\noindent\textbf{Your Answer:} \fbox{\rule{0pt}{1.3ex}\hspace{3cm}}

\tcblower

\textbf{Official Key:} \textbf{\large 4 to 4} \hfill \textit{\small [Medium]}
\vspace{2pt}\hrule height 0.3pt\vspace{3pt}
\textbf{Topic:} Programming and Data Structures \\
\textbf{Subtopic:} \textit{Trees, binary search trees, binary heaps}
\vspace{2pt}\hrule height 0.3pt\vspace{3pt}
\textbf{Guide / Hint:} \\
{\small In postorder traversal, the root is the last element (1). Splitting inorder $[8,6,9,4,7,2,5]$ and $[3]$ around root 1 gives left subtree nodes and right subtree node 3. Recursively constructing the unique binary tree reveals the longest leaf path: $1 \to 2 \to 4 \to 6 \to 8$ (or 9), which has 4 edges. Thus, the height in terms of edges is 4.}

\end{tcolorbox}
\vspace{6pt}\hrule height 0.35pt\vspace{6pt}
\noindent
\begin{tcolorbox}[
    blanker,
    sidebyside,
    sidebyside align=top,
    lefthand width=0.55\textwidth,
    top=4pt, bottom=4pt, left=0pt, right=0pt,
    middle=16pt,
    underlay={
        \draw[dashed, black!80, line width=0.8pt, dash pattern=on 4pt off 3pt]
            ([xshift=0.55\textwidth+8pt]frame.north west) --
            ([xshift=0.55\textwidth+8pt]frame.south west);
    }
]
\textbf{Q.3} \quad {\footnotesize\textsc{[NAT $\cdot$ 1 Mark(s) $\cdot$ No Neg]}} \hfill {\footnotesize\textbf{GATE 2018 CS} \quad \tiny\texttt{2018\_CS\_GA\_Q21}}
\vspace{2pt}\hrule height 0.3pt\vspace{4pt}

Consider the following C program:
\begin{lstlisting}[language=C]
#include <stdio.h>
int counter = 0;
\par\medskip
int calc(int a, int b) {
    int c;
    counter++;
    if (b <= 3) return (a + b);
    else {
        c = b / 3;
        return (calc(a, c));
    }
}
\par\medskip
int main() {
    calc(4, 81);
    printf("%d", counter);
    return 0;
}
\end{lstlisting}
The output of this program is \_\_\_\_\_\_\_\_.

\par\vspace{6pt}\noindent\textbf{Your Answer:} \fbox{\rule{0pt}{1.3ex}\hspace{3cm}}

\tcblower

\textbf{Official Key:} \textbf{\large 4 to 4} \hfill \textit{\small [Medium]}
\vspace{2pt}\hrule height 0.3pt\vspace{3pt}
\textbf{Topic:} Programming and Data Structures \\
\textbf{Subtopic:} \textit{Recursion}
\vspace{2pt}\hrule height 0.3pt\vspace{3pt}
\textbf{Guide / Hint:} \\
{\small Function `calc(a, b)` computes $\log_3(b)$ recursively: `calc(4, 81)` calls `calc(4, 27)` $\to$ `calc(4, 9)` $\to$ `calc(4, 3)`. At $b = 3$, `b <= 3` is satisfied and it returns. The variable `counter` is incremented once per invocation, resulting in exactly 4 calls.}

\end{tcolorbox}
\vspace{6pt}\hrule height 0.35pt\vspace{6pt}
\noindent
\begin{tcolorbox}[
    blanker,
    sidebyside,
    sidebyside align=top,
    lefthand width=0.55\textwidth,
    top=4pt, bottom=4pt, left=0pt, right=0pt,
    middle=16pt,
    underlay={
        \draw[dashed, black!80, line width=0.8pt, dash pattern=on 4pt off 3pt]
            ([xshift=0.55\textwidth+8pt]frame.north west) --
            ([xshift=0.55\textwidth+8pt]frame.south west);
    }
]
\textbf{Q.4} \quad {\footnotesize\textsc{[MCQ $\cdot$ 2 Mark(s) $\cdot$ Neg: 0.66]}} \hfill {\footnotesize\textbf{GATE 2018 CS} \quad \tiny\texttt{2018\_CS\_GA\_Q29}}
\vspace{2pt}\hrule height 0.3pt\vspace{4pt}

Consider the following C program:
\begin{lstlisting}[language=C]
#include <stdio.h>
\par\medskip
void fun1(char *s1, char *s2) {
    char *tmp;
    tmp = s1;
    s1 = s2;
    s2 = tmp;
}
\par\medskip
void fun2(char **s1, char **s2) {
    char *tmp;
    tmp = *s1;
    *s1 = *s2;
    *s2 = tmp;
}
\par\medskip
int main() {
    char *str1 = "Hi", *str2 = "Bye";
    fun1(str1, str2);
    printf("%s %s ", str1, str2);
    fun2(&str1, &str2);
    printf("%s %s", str1, str2);
    return 0;
}
\end{lstlisting}
The output of the program is

\par\vspace{6pt}\noindent
\begin{tabular}{@{}p{0.48\linewidth}@{\hspace{0.04\linewidth}}p{0.48\linewidth}@{}}
$\square$\ \textbf{(A)} Hi Bye Bye Hi & $\square$\ \textbf{(B)} Hi Bye Hi Bye \\
\addlinespace[3pt]
$\square$\ \textbf{(C)} Bye Hi Bye Hi & $\square$\ \textbf{(D)} Bye Hi Hi Bye
\end{tabular}

\tcblower

\textbf{Official Key:} \textbf{\large A} \hfill \textit{\small [Medium]}
\vspace{2pt}\hrule height 0.3pt\vspace{3pt}
\textbf{Topic:} Programming and Data Structures \\
\textbf{Subtopic:} \textit{Programming in C}
\vspace{2pt}\hrule height 0.3pt\vspace{3pt}
\textbf{Guide / Hint:} \\
{\small In `fun1`, pointers `s1` and `s2` are passed by value, so the pointer swap only alters local variables and has no effect on `str1` and `str2` in `main` (printing "Hi Bye"). In `fun2`, the addresses of the pointers (`char**`) are passed, so dereferencing swaps the actual pointer values `str1` and `str2` (printing "Bye Hi").}

\end{tcolorbox}
\vspace{6pt}\hrule height 0.35pt\vspace{6pt}
\noindent
\begin{tcolorbox}[
    blanker,
    sidebyside,
    sidebyside align=top,
    lefthand width=0.55\textwidth,
    top=4pt, bottom=4pt, left=0pt, right=0pt,
    middle=16pt,
    underlay={
        \draw[dashed, black!80, line width=0.8pt, dash pattern=on 4pt off 3pt]
            ([xshift=0.55\textwidth+8pt]frame.north west) --
            ([xshift=0.55\textwidth+8pt]frame.south west);
    }
]
\textbf{Q.5} \quad {\footnotesize\textsc{[NAT $\cdot$ 2 Mark(s) $\cdot$ No Neg]}} \hfill {\footnotesize\textbf{GATE 2018 CS} \quad \tiny\texttt{2018\_CS\_GA\_Q45}}
\vspace{2pt}\hrule height 0.3pt\vspace{4pt}

Consider the following pseudo-code function where $x$ and $y$ are positive integers:
\begin{lstlisting}[language=C]
Count(x, y) {
    if (y != 1) {
        if (x != 1) {
            print("*");
            Count(x / 2, y);
        } else {
            Count(1024, y - 1);
        }
    }
}
\end{lstlisting}
The total number of times `print("*")` is executed for `Count(1024, 1024)` is \_\_\_\_\_\_\_\_.

\par\vspace{6pt}\noindent\textbf{Your Answer:} \fbox{\rule{0pt}{1.3ex}\hspace{3cm}}

\tcblower

\textbf{Official Key:} \textbf{\large 10230 to 10230} \hfill \textit{\small [Medium]}
\vspace{2pt}\hrule height 0.3pt\vspace{3pt}
\textbf{Topic:} Programming and Data Structures \\
\textbf{Subtopic:} \textit{Recursion}
\vspace{2pt}\hrule height 0.3pt\vspace{3pt}
\textbf{Guide / Hint:} \\
{\small In `Count(x, y)`, outer recursion on $y$ runs for $y = 1024, 1023, \dots, 2$ ($1023$ values). For each $y$, the inner branch halves $x$ until $x = 1$: starting from $x = 1024$, this loop prints 10 times (since $\log_2(1024) = 10$). Total prints $= 1023 \times 10 = 10230$.}

\end{tcolorbox}
\vspace{6pt}\hrule height 0.35pt\vspace{6pt}
\noindent
\begin{tcolorbox}[
    blanker,
    sidebyside,
    sidebyside align=top,
    lefthand width=0.55\textwidth,
    top=4pt, bottom=4pt, left=0pt, right=0pt,
    middle=16pt,
    underlay={
        \draw[dashed, black!80, line width=0.8pt, dash pattern=on 4pt off 3pt]
            ([xshift=0.55\textwidth+8pt]frame.north west) --
            ([xshift=0.55\textwidth+8pt]frame.south west);
    }
]
\textbf{Q.6} \quad {\footnotesize\textsc{[NAT $\cdot$ 2 Mark(s) $\cdot$ No Neg]}} \hfill {\footnotesize\textbf{GATE 2018 CS} \quad \tiny\texttt{2018\_CS\_GA\_Q46}}
\vspace{2pt}\hrule height 0.3pt\vspace{4pt}

The number of possible min-heaps containing each value from $\{1, 2, 3, 4, 5, 6, 7\}$ exactly once is \_\_\_\_\_\_\_\_.

\par\vspace{6pt}\noindent\textbf{Your Answer:} \fbox{\rule{0pt}{1.3ex}\hspace{3cm}}

\tcblower

\textbf{Official Key:} \textbf{\large 80 to 80} \hfill \textit{\small [Medium]}
\vspace{2pt}\hrule height 0.3pt\vspace{3pt}
\textbf{Topic:} Programming and Data Structures \\
\textbf{Subtopic:} \textit{Trees, binary search trees, binary heaps}
\vspace{2pt}\hrule height 0.3pt\vspace{3pt}
\textbf{Guide / Hint:} \\
{\small In a min-heap of 7 elements, the root must contain the minimum element (1). The left subtree has size 3 and the right subtree has size 3. Choosing 3 elements for the left subtree from the remaining 6 elements can be done in $\binom{6}{3} = 20$ ways. A 3-element min-heap has $\binom{2}{1} = 2$ valid configurations. Hence, total valid min-heaps $= 20 \times 2 \times 2 = 80$.}

\end{tcolorbox}
\vspace{6pt}\hrule height 0.35pt\vspace{6pt}
\noindent
\begin{tcolorbox}[
    blanker,
    sidebyside,
    sidebyside align=top,
    lefthand width=0.55\textwidth,
    top=4pt, bottom=4pt, left=0pt, right=0pt,
    middle=16pt,
    underlay={
        \draw[dashed, black!80, line width=0.8pt, dash pattern=on 4pt off 3pt]
            ([xshift=0.55\textwidth+8pt]frame.north west) --
            ([xshift=0.55\textwidth+8pt]frame.south west);
    }
]
\textbf{Q.7} \quad {\footnotesize\textsc{[MCQ $\cdot$ 1 Mark(s) $\cdot$ Neg: 0.33]}} \hfill {\footnotesize\textbf{GATE 2022 CS} \quad \tiny\texttt{2022\_CS\_CS\_Q15}}
\vspace{2pt}\hrule height 0.3pt\vspace{4pt}

Consider the problem of reversing a singly linked list. Which one of the following statements is TRUE about the time complexity of algorithms that solve the above problem in $\mathcal{O}(1)$ space?

\par\vspace{6pt}\noindent
\begin{tabular}{@{}p{0.98\linewidth}@{}}
$\square$\ \textbf{(A)} The best algorithm for the problem takes $\Theta(n)$ time in the worst case. \\
\addlinespace[3pt]
$\square$\ \textbf{(B)} The best algorithm for the problem takes $\Theta(n \log n)$ time in the worst case. \\
\addlinespace[3pt]
$\square$\ \textbf{(C)} The best algorithm for the problem takes $\Theta(n^2)$ time in the worst case. \\
\addlinespace[3pt]
$\square$\ \textbf{(D)} It is not possible to reverse a singly linked list in $\mathcal{O}(1)$ space.
\end{tabular}

\tcblower

\textbf{Official Key:} \textbf{\large A} \hfill \textit{\small [Easy]}
\vspace{2pt}\hrule height 0.3pt\vspace{3pt}
\textbf{Topic:} Programming and Data Structures \\
\textbf{Subtopic:} \textit{Arrays, stacks, queues, linked lists}
\vspace{2pt}\hrule height 0.3pt\vspace{3pt}
\textbf{Guide / Hint:} \\
{\small A singly linked list can be reversed in-place in $\mathcal{O}(1)$ auxiliary space using three moving pointers (`prev`, `curr`, `next`) by reversing each pointer sequentially in a single pass of $n$ nodes, achieving optimal $\Theta(n)$ worst-case time.}

\end{tcolorbox}
\vspace{6pt}\hrule height 0.35pt\vspace{6pt}
\noindent
\begin{tcolorbox}[
    blanker,
    sidebyside,
    sidebyside align=top,
    lefthand width=0.55\textwidth,
    top=4pt, bottom=4pt, left=0pt, right=0pt,
    middle=16pt,
    underlay={
        \draw[dashed, black!80, line width=0.8pt, dash pattern=on 4pt off 3pt]
            ([xshift=0.55\textwidth+8pt]frame.north west) --
            ([xshift=0.55\textwidth+8pt]frame.south west);
    }
]
\textbf{Q.8} \quad {\footnotesize\textsc{[MCQ $\cdot$ 1 Mark(s) $\cdot$ Neg: 0.33]}} \hfill {\footnotesize\textbf{GATE 2022 CS} \quad \tiny\texttt{2022\_CS\_CS\_Q21}}
\vspace{2pt}\hrule height 0.3pt\vspace{4pt}

What is printed by the following ANSI C program?
\begin{lstlisting}[language=C]
#include <stdio.h>
int main(int argc, char *argv[]) {
    int x = 1, z[2] = {10, 11};
    int *p = NULL;
    p = &x;
    *p = 10;
    p = &z[1];
    *(&z[0] + 1) += 3;
    printf("%d, %d, %d\n", x, z[0], z[1]);
    return 0;
}
\end{lstlisting}

\par\vspace{6pt}\noindent
\begin{tabular}{@{}p{0.48\linewidth}@{\hspace{0.04\linewidth}}p{0.48\linewidth}@{}}
$\square$\ \textbf{(A)} 1, 10, 11 & $\square$\ \textbf{(B)} 1, 10, 14 \\
\addlinespace[3pt]
$\square$\ \textbf{(C)} 10, 14, 11 & $\square$\ \textbf{(D)} 10, 10, 14
\end{tabular}

\tcblower

\textbf{Official Key:} \textbf{\large D} \hfill \textit{\small [Easy]}
\vspace{2pt}\hrule height 0.3pt\vspace{3pt}
\textbf{Topic:} Programming and Data Structures \\
\textbf{Subtopic:} \textit{Programming in C}
\vspace{2pt}\hrule height 0.3pt\vspace{3pt}
\textbf{Guide / Hint:} \\
{\small Initially $x = 1, z[0] = 10, z[1] = 11$. The assignment `p = \&x; *p = 10;` sets $x = 10$. The line `*(\&z[0] + 1) += 3;` targets `z[1]`, updating $z[1] = 11 + 3 = 14$. The array element $z[0]$ remains 10. The output printed is `10, 10, 14`.}

\end{tcolorbox}
\vspace{6pt}\hrule height 0.35pt\vspace{6pt}
\noindent
\begin{tcolorbox}[
    blanker,
    sidebyside,
    sidebyside align=top,
    lefthand width=0.55\textwidth,
    top=4pt, bottom=4pt, left=0pt, right=0pt,
    middle=16pt,
    underlay={
        \draw[dashed, black!80, line width=0.8pt, dash pattern=on 4pt off 3pt]
            ([xshift=0.55\textwidth+8pt]frame.north west) --
            ([xshift=0.55\textwidth+8pt]frame.south west);
    }
]
\textbf{Q.9} \quad {\footnotesize\textsc{[NAT $\cdot$ 1 Mark(s) $\cdot$ No Neg]}} \hfill {\footnotesize\textbf{GATE 2022 CS} \quad \tiny\texttt{2022\_CS\_CS\_Q28}}
\vspace{2pt}\hrule height 0.3pt\vspace{4pt}

Suppose a binary search tree with 1000 distinct elements is also a complete binary tree. The tree is stored using the array representation of binary heap trees. Assuming that the array indices start with 0, the 3rd largest element of the tree is stored at index \_\_\_\_\_\_\_\_\_\_\_\_\. (Answer in integer)

\par\vspace{6pt}\noindent\textbf{Your Answer:} \fbox{\rule{0pt}{1.3ex}\hspace{3cm}}

\tcblower

\textbf{Official Key:} \textbf{\large 509 to 509} \hfill \textit{\small [Hard]}
\vspace{2pt}\hrule height 0.3pt\vspace{3pt}
\textbf{Topic:} Programming and Data Structures \\
\textbf{Subtopic:} \textit{Trees, binary search trees, binary heaps}
\vspace{2pt}\hrule height 0.3pt\vspace{3pt}
\textbf{Guide / Hint:} \\
{\small In a complete binary tree BST with 1000 nodes (indices 0 to 999), the rightmost branch contains nodes 0, 2, 6, 14, 30, 62, 126, 254, 510. Node 510 is a leaf and holds the maximum element (1st largest). Node 254 is its parent and holds the 2nd largest element. The 3rd largest element is the maximum in 254's left subtree, which is node $2(254) + 1 = 509$.}

\end{tcolorbox}
\vspace{6pt}\hrule height 0.35pt\vspace{6pt}
\noindent
\begin{tcolorbox}[
    blanker,
    sidebyside,
    sidebyside align=top,
    lefthand width=0.55\textwidth,
    top=4pt, bottom=4pt, left=0pt, right=0pt,
    middle=16pt,
    underlay={
        \draw[dashed, black!80, line width=0.8pt, dash pattern=on 4pt off 3pt]
            ([xshift=0.55\textwidth+8pt]frame.north west) --
            ([xshift=0.55\textwidth+8pt]frame.south west);
    }
]
\textbf{Q.10} \quad {\footnotesize\textsc{[MCQ $\cdot$ 2 Mark(s) $\cdot$ Neg: 0.66]}} \hfill {\footnotesize\textbf{GATE 2022 CS} \quad \tiny\texttt{2022\_CS\_CS\_Q43}}
\vspace{2pt}\hrule height 0.3pt\vspace{4pt}

What is printed by the following ANSI C program?
\begin{lstlisting}[language=C]
#include <stdio.h>
int main(int argc, char *argv[]) {
    int a[3][3][3] = {
        {1, 2, 3, 4, 5, 6, 7, 8, 9},
        {10, 11, 12, 13, 14, 15, 16, 17, 18},
        {19, 20, 21, 22, 23, 24, 25, 26, 27}
    };
    int i = 0, j = 0, k = 0;
    for(i = 0; i < 3; i++) {
        for(k = 0; k < 3; k++)
            printf("%d ", a[i][j][k]);
        printf("\n");
    }
    return 0;
}
\end{lstlisting}

\par\vspace{6pt}\noindent
\begin{tabular}{@{}p{0.48\linewidth}@{\hspace{0.04\linewidth}}p{0.48\linewidth}@{}}
$\square$\ \textbf{(A)} 1 2 3 \\ 10 11 12 \\ 19 20 21 & $\square$\ \textbf{(B)} 1 4 7 \\ 10 13 16 \\ 19 22 25 \\
\addlinespace[3pt]
$\square$\ \textbf{(C)} 1 2 3 \\ 4 5 6 \\ 7 8 9 & $\square$\ \textbf{(D)} 1 2 3 \\ 13 14 15 \\ 25 26 27
\end{tabular}

\tcblower

\textbf{Official Key:} \textbf{\large A} \hfill \textit{\small [Easy]}
\vspace{2pt}\hrule height 0.3pt\vspace{3pt}
\textbf{Topic:} Programming and Data Structures \\
\textbf{Subtopic:} \textit{Programming in C}
\vspace{2pt}\hrule height 0.3pt\vspace{3pt}
\textbf{Guide / Hint:} \\
{\small In the nested loops, index $j = 0$ is fixed throughout. For $i = 0$, $k$ iterates from 0 to 2, printing `a[0][0][0], a[0][0][1], a[0][0][2]` ($1\ 2\ 3$). For $i = 1$, it prints `a[1][0][0], a[1][0][1], a[1][0][2]` ($10\ 11\ 12$). For $i = 2$, it prints `a[2][0][0], a[2][0][1], a[2][0][2]` ($19\ 20\ 21$).}

\end{tcolorbox}
\vspace{6pt}\hrule height 0.35pt\vspace{6pt}
\noindent
\begin{tcolorbox}[
    blanker,
    sidebyside,
    sidebyside align=top,
    lefthand width=0.55\textwidth,
    top=4pt, bottom=4pt, left=0pt, right=0pt,
    middle=16pt,
    underlay={
        \draw[dashed, black!80, line width=0.8pt, dash pattern=on 4pt off 3pt]
            ([xshift=0.55\textwidth+8pt]frame.north west) --
            ([xshift=0.55\textwidth+8pt]frame.south west);
    }
]
\textbf{Q.11} \quad {\footnotesize\textsc{[MCQ $\cdot$ 2 Mark(s) $\cdot$ Neg: 0.66]}} \hfill {\footnotesize\textbf{GATE 2022 CS} \quad \tiny\texttt{2022\_CS\_CS\_Q44}}
\vspace{2pt}\hrule height 0.3pt\vspace{4pt}

What is printed by the following ANSI C program?
\begin{lstlisting}[language=C]
#include <stdio.h>
int main(int argc, char *argv[]) {
    char a = 'P';
    char b = 'x';
    char c = (a & b) + '*';
    char d = (a | b) - '-';
    char e = (a ^ b) + '+';
    printf("%c %c %c\n", c, d, e);
    return 0;
}
\end{lstlisting}

\par\vspace{6pt}\noindent
\begin{tabular}{@{}p{0.48\linewidth}@{\hspace{0.04\linewidth}}p{0.48\linewidth}@{}}
$\square$\ \textbf{(A)} z K S & $\square$\ \textbf{(B)} 122 75 83 \\
\addlinespace[3pt]
$\square$\ \textbf{(C)} * - + & $\square$\ \textbf{(D)} P x +
\end{tabular}

\tcblower

\textbf{Official Key:} \textbf{\large A} \hfill \textit{\small [Medium]}
\vspace{2pt}\hrule height 0.3pt\vspace{3pt}
\textbf{Topic:} Programming and Data Structures \\
\textbf{Subtopic:} \textit{Programming in C}
\vspace{2pt}\hrule height 0.3pt\vspace{3pt}
\textbf{Guide / Hint:} \\
{\small Character values: `a = 'P'` (80), `b = 'x'` (120). $c = (80 \& 120) + '*' = 80 + 42 = 122$, which is ASCII for 'z'. $d = (80 \mid 120) - '-' = 120 - 45 = 75$, which is ASCII for 'K'. $e = (80 \oplus 120) + '+' = 40 + 43 = 83$, which is ASCII for 'S'. The printed characters are `z K S`.}

\end{tcolorbox}
\vspace{6pt}\hrule height 0.35pt\vspace{6pt}
\noindent
\begin{tcolorbox}[
    blanker,
    sidebyside,
    sidebyside align=top,
    lefthand width=0.55\textwidth,
    top=4pt, bottom=4pt, left=0pt, right=0pt,
    middle=16pt,
    underlay={
        \draw[dashed, black!80, line width=0.8pt, dash pattern=on 4pt off 3pt]
            ([xshift=0.55\textwidth+8pt]frame.north west) --
            ([xshift=0.55\textwidth+8pt]frame.south west);
    }
]
\textbf{Q.12} \quad {\footnotesize\textsc{[NAT $\cdot$ 2 Mark(s) $\cdot$ No Neg]}} \hfill {\footnotesize\textbf{GATE 2022 CS} \quad \tiny\texttt{2022\_CS\_CS\_Q62}}
\vspace{2pt}\hrule height 0.3pt\vspace{4pt}

Consider the queues $Q_1$ containing four elements and $Q_2$ containing none. The only operations allowed on these two queues are `Enqueue(Q, element)` and `Dequeue(Q)`. The minimum number of `Enqueue` operations on $Q_1$ required to place the elements of $Q_1$ in $Q_2$ in reverse order without using any additional storage is \_\_\_\_\_\_\_\_\_\_\_\. (Answer in integer)

\par\vspace{6pt}\noindent\textbf{Your Answer:} \fbox{\rule{0pt}{1.3ex}\hspace{3cm}}

\tcblower

\textbf{Official Key:} \textbf{\large 0 to 0} \hfill \textit{\small [Medium]}
\vspace{2pt}\hrule height 0.3pt\vspace{3pt}
\textbf{Topic:} Programming and Data Structures \\
\textbf{Subtopic:} \textit{Arrays, stacks, queues, linked lists}
\vspace{2pt}\hrule height 0.3pt\vspace{3pt}
\textbf{Guide / Hint:} \\
{\small To place the 4 elements of $Q_1$ into $Q_2$ in reverse order, elements are dequeued one by one from $Q_1$ and placed into $Q_2$. Internal order manipulation can be performed entirely by cycling elements within $Q_2$ via dequeue/enqueue on $Q_2$, requiring 0 enqueue operations on $Q_1$.}

\end{tcolorbox}
\vspace{6pt}\hrule height 0.35pt\vspace{6pt}
\noindent
\begin{tcolorbox}[
    blanker,
    sidebyside,
    sidebyside align=top,
    lefthand width=0.55\textwidth,
    top=4pt, bottom=4pt, left=0pt, right=0pt,
    middle=16pt,
    underlay={
        \draw[dashed, black!80, line width=0.8pt, dash pattern=on 4pt off 3pt]
            ([xshift=0.55\textwidth+8pt]frame.north west) --
            ([xshift=0.55\textwidth+8pt]frame.south west);
    }
]
\textbf{Q.13} \quad {\footnotesize\textsc{[MCQ $\cdot$ 1 Mark(s) $\cdot$ Neg: 0.33]}} \hfill {\footnotesize\textbf{GATE 2023 CS} \quad \tiny\texttt{2023\_CS\_CS\_Q12}}
\vspace{2pt}\hrule height 0.3pt\vspace{4pt}

Which one of the following sequences when stored in an array at locations $A[1], \dots, A[10]$ forms a max-heap?

\par\vspace{6pt}\noindent
\begin{tabular}{@{}p{0.98\linewidth}@{}}
$\square$\ \textbf{(A)} 23, 17, 10, 6, 13, 14, 1, 5, 7, 12 \\
\addlinespace[3pt]
$\square$\ \textbf{(B)} 23, 17, 14, 7, 13, 10, 1, 5, 6, 12 \\
\addlinespace[3pt]
$\square$\ \textbf{(C)} 23, 17, 14, 6, 13, 10, 1, 5, 7, 15 \\
\addlinespace[3pt]
$\square$\ \textbf{(D)} 23, 14, 17, 1, 10, 13, 16, 12, 7, 5
\end{tabular}

\tcblower

\textbf{Official Key:} \textbf{\large B} \hfill \textit{\small [Easy]}
\vspace{2pt}\hrule height 0.3pt\vspace{3pt}
\textbf{Topic:} Programming and Data Structures \\
\textbf{Subtopic:} \textit{Trees, binary search trees, binary heaps}
\vspace{2pt}\hrule height 0.3pt\vspace{3pt}
\textbf{Guide / Hint:} \\
{\small In a 1-indexed max-heap, every node at index $i$ must satisfy $A[i] \ge A[2i]$ and $A[i] \ge A[2i+1]$. For sequence B $[23, 17, 14, 7, 13, 10, 1, 5, 6, 12]$: $A[1]=23 \ge A[2]=17, A[3]=14$; $A[2]=17 \ge A[4]=7, A[5]=13$; $A[3]=14 \ge A[6]=10, A[7]=1$; $A[4]=7 \ge A[8]=5, A[9]=6$; $A[5]=13 \ge A[10]=12$. Every parent is greater than or equal to its children, forming a valid max-heap.}

\end{tcolorbox}
\vspace{6pt}\hrule height 0.35pt\vspace{6pt}
\noindent
\begin{tcolorbox}[
    blanker,
    sidebyside,
    sidebyside align=top,
    lefthand width=0.55\textwidth,
    top=4pt, bottom=4pt, left=0pt, right=0pt,
    middle=16pt,
    underlay={
        \draw[dashed, black!80, line width=0.8pt, dash pattern=on 4pt off 3pt]
            ([xshift=0.55\textwidth+8pt]frame.north west) --
            ([xshift=0.55\textwidth+8pt]frame.south west);
    }
]
\textbf{Q.14} \quad {\footnotesize\textsc{[MCQ $\cdot$ 1 Mark(s) $\cdot$ Neg: 0.33]}} \hfill {\footnotesize\textbf{GATE 2023 CS} \quad \tiny\texttt{2023\_CS\_CS\_Q13}}
\vspace{2pt}\hrule height 0.3pt\vspace{4pt}

Let `SLLdel` be a function that deletes a node in a singly-linked list given a pointer to the node and a pointer to the head of the list. Similarly, let `DLLdel` be another function that deletes a node in a doubly-linked list given a pointer to the node and a pointer to the head of the list.
Let $n$ denote the number of nodes in each of the linked lists. Which one of the following choices is TRUE about the worst-case time complexity of `SLLdel` and `DLLdel`?

\par\vspace{6pt}\noindent
\begin{tabular}{@{}p{0.98\linewidth}@{}}
$\square$\ \textbf{(A)} `SLLdel` is $\mathcal{O}(1)$ and `DLLdel` is $\mathcal{O}(n)$ \\
\addlinespace[3pt]
$\square$\ \textbf{(B)} Both `SLLdel` and `DLLdel` are $\mathcal{O}(\log n)$ \\
\addlinespace[3pt]
$\square$\ \textbf{(C)} Both `SLLdel` and `DLLdel` are $\mathcal{O}(1)$ \\
\addlinespace[3pt]
$\square$\ \textbf{(D)} `SLLdel` is $\mathcal{O}(n)$ and `DLLdel` is $\mathcal{O}(1)$
\end{tabular}

\tcblower

\textbf{Official Key:} \textbf{\large D} \hfill \textit{\small [Easy]}
\vspace{2pt}\hrule height 0.3pt\vspace{3pt}
\textbf{Topic:} Programming and Data Structures \\
\textbf{Subtopic:} \textit{Arrays, stacks, queues, linked lists}
\vspace{2pt}\hrule height 0.3pt\vspace{3pt}
\textbf{Guide / Hint:} \\
{\small In a singly-linked list, deleting a given node requires finding its predecessor node to bypass it, requiring an $\mathcal{O}(n)$ traversal from the head in the worst case. In a doubly-linked list, each node directly points to its predecessor (`node->prev`), allowing immediate pointer adjustments in $\mathcal{O}(1)$ time.}

\end{tcolorbox}
\vspace{6pt}\hrule height 0.35pt\vspace{6pt}
\noindent
\begin{tcolorbox}[
    blanker,
    sidebyside,
    sidebyside align=top,
    lefthand width=0.55\textwidth,
    top=4pt, bottom=4pt, left=0pt, right=0pt,
    middle=16pt,
    underlay={
        \draw[dashed, black!80, line width=0.8pt, dash pattern=on 4pt off 3pt]
            ([xshift=0.55\textwidth+8pt]frame.north west) --
            ([xshift=0.55\textwidth+8pt]frame.south west);
    }
]
\textbf{Q.15} \quad {\footnotesize\textsc{[NAT $\cdot$ 1 Mark(s) $\cdot$ No Neg]}} \hfill {\footnotesize\textbf{GATE 2023 CS} \quad \tiny\texttt{2023\_CS\_CS\_Q35}}
\vspace{2pt}\hrule height 0.3pt\vspace{4pt}

The integer value printed by the ANSI-C program given below is \_\_\_\_\_\_\_\_\_\_\. (Answer in integer)
\begin{lstlisting}[language=C]
#include <stdio.h>
int funcp() {
    static int x = 1;
    x++;
    return x;
}
int main() {
    int x, y;
    x = funcp();
    y = funcp() + x;
    printf("%d\n", (x + y));
    return 0;
}
\end{lstlisting}

\par\vspace{6pt}\noindent\textbf{Your Answer:} \fbox{\rule{0pt}{1.3ex}\hspace{3cm}}

\tcblower

\textbf{Official Key:} \textbf{\large 7 to 7} \hfill \textit{\small [Easy]}
\vspace{2pt}\hrule height 0.3pt\vspace{3pt}
\textbf{Topic:} Programming and Data Structures \\
\textbf{Subtopic:} \textit{Programming in C}
\vspace{2pt}\hrule height 0.3pt\vspace{3pt}
\textbf{Guide / Hint:} \\
{\small In `funcp()`, the static variable $x$ persists across invocations. On first call `x = funcp()`, static $x$ increments from 1 to 2 and returns 2, setting local $x = 2$. In `y = funcp() + x`, `funcp()` increments static $x$ from 2 to 3 and returns 3, so $y = 3 + 2 = 5$. Finally, `printf` prints $x + y = 2 + 5 = 7$.}

\end{tcolorbox}
\vspace{6pt}\hrule height 0.35pt\vspace{6pt}
\noindent
\begin{tcolorbox}[
    blanker,
    sidebyside,
    sidebyside align=top,
    lefthand width=0.55\textwidth,
    top=4pt, bottom=4pt, left=0pt, right=0pt,
    middle=16pt,
    underlay={
        \draw[dashed, black!80, line width=0.8pt, dash pattern=on 4pt off 3pt]
            ([xshift=0.55\textwidth+8pt]frame.north west) --
            ([xshift=0.55\textwidth+8pt]frame.south west);
    }
]
\textbf{Q.16} \quad {\footnotesize\textsc{[MCQ $\cdot$ 2 Mark(s) $\cdot$ Neg: 0.66]}} \hfill {\footnotesize\textbf{GATE 2023 CS} \quad \tiny\texttt{2023\_CS\_CS\_Q46}}
\vspace{2pt}\hrule height 0.3pt\vspace{4pt}

Let $A$ be a priority queue for maintaining a set of elements. Suppose $A$ is implemented using a max-heap data structure. The operation $\text{Extract-Max}(A)$ extracts and deletes the maximum element from $A$. The operation $\text{Insert}(A, \text{key})$ inserts a new element $\text{key}$ in $A$. The properties of a max-heap are preserved at the end of each of these operations.
When $A$ contains $n$ elements, which one of the following statements about the worst case running time of these two operations is TRUE?

\par\vspace{6pt}\noindent
\begin{tabular}{@{}p{0.98\linewidth}@{}}
$\square$\ \textbf{(A)} Both `Extract-Max(A)` and `Insert(A,key)` run in $\mathcal{O}(1)$. \\
\addlinespace[3pt]
$\square$\ \textbf{(B)} Both `Extract-Max(A)` and `Insert(A,key)` run in $\mathcal{O}(\log n)$. \\
\addlinespace[3pt]
$\square$\ \textbf{(C)} `Extract-Max(A)` runs in $\mathcal{O}(1)$ whereas `Insert(A,key)` runs in $\mathcal{O}(n)$. \\
\addlinespace[3pt]
$\square$\ \textbf{(D)} `Extract-Max(A)` runs in $\mathcal{O}(1)$ whereas `Insert(A,key)` runs in $\mathcal{O}(\log n)$.
\end{tabular}

\tcblower

\textbf{Official Key:} \textbf{\large B} \hfill \textit{\small [Easy]}
\vspace{2pt}\hrule height 0.3pt\vspace{3pt}
\textbf{Topic:} Programming and Data Structures \\
\textbf{Subtopic:} \textit{Trees, binary search trees, binary heaps}
\vspace{2pt}\hrule height 0.3pt\vspace{3pt}
\textbf{Guide / Hint:} \\
{\small In a binary max-heap of $n$ elements, `Extract-Max` replaces the root with the last element and performs `maxHeapify` downward along the height $\lfloor \log_2 n \rfloor$, running in $\mathcal{O}(\log n)$ worst-case time. `Insert` places the new key at the next leaf position and sifts upward along the height, also running in $\mathcal{O}(\log n)$ worst-case time.}

\end{tcolorbox}
\vspace{6pt}\hrule height 0.35pt\vspace{6pt}
\noindent
\begin{tcolorbox}[
    blanker,
    sidebyside,
    sidebyside align=top,
    lefthand width=0.55\textwidth,
    top=4pt, bottom=4pt, left=0pt, right=0pt,
    middle=16pt,
    underlay={
        \draw[dashed, black!80, line width=0.8pt, dash pattern=on 4pt off 3pt]
            ([xshift=0.55\textwidth+8pt]frame.north west) --
            ([xshift=0.55\textwidth+8pt]frame.south west);
    }
]
\textbf{Q.17} \quad {\footnotesize\textsc{[MCQ $\cdot$ 2 Mark(s) $\cdot$ Neg: 0.66]}} \hfill {\footnotesize\textbf{GATE 2023 CS} \quad \tiny\texttt{2023\_CS\_CS\_Q47}}
\vspace{2pt}\hrule height 0.3pt\vspace{4pt}

Consider the C function `foo` and the binary tree shown:
\begin{lstlisting}[language=C]
typedef struct node {
    int val;
    struct node *left, *right;
} node;
int foo(node *p) {
    int retval;
    if (p == NULL)
        return 0;
    else {
        retval = p->val + foo(p->left) + foo(p->right);
        printf("%d ", retval);
        return retval;
    }
}
\end{lstlisting}
In the binary tree, root is 10; its left child is 5 with children 3 and 8; its right child is 11 with right child 13.
When `foo` is called with a pointer to the root node of the given binary tree, what will it print?

\par\vspace{6pt}\noindent
\begin{tabular}{@{}p{0.48\linewidth}@{\hspace{0.04\linewidth}}p{0.48\linewidth}@{}}
$\square$\ \textbf{(A)} 3 8 5 13 11 10 & $\square$\ \textbf{(B)} 3 5 8 10 11 13 \\
\addlinespace[3pt]
$\square$\ \textbf{(C)} 3 8 16 13 24 50 & $\square$\ \textbf{(D)} 3 16 8 50 24 13
\end{tabular}

\tcblower

\textbf{Official Key:} \textbf{\large C} \hfill \textit{\small [Medium]}
\vspace{2pt}\hrule height 0.3pt\vspace{3pt}
\textbf{Topic:} Programming and Data Structures \\
\textbf{Subtopic:} \textit{Recursion}
\vspace{2pt}\hrule height 0.3pt\vspace{3pt}
\textbf{Guide / Hint:} \\
{\small Function `foo` performs post-order traversal, recursively summing `p->val + foo(p->left) + foo(p->right)` and printing each subtree sum. Leaf 3 prints 3; leaf 8 prints 8; node 5 sums $5 + 3 + 8 = 16$ (prints 16); leaf 13 prints 13; node 11 sums $11 + 0 + 13 = 24$ (prints 24); root 10 sums $10 + 16 + 24 = 50$ (prints 50). The printed output sequence is 3 8 16 13 24 50.}

\end{tcolorbox}
\vspace{6pt}\hrule height 0.35pt\vspace{6pt}
\noindent
\begin{tcolorbox}[
    blanker,
    sidebyside,
    sidebyside align=top,
    lefthand width=0.55\textwidth,
    top=4pt, bottom=4pt, left=0pt, right=0pt,
    middle=16pt,
    underlay={
        \draw[dashed, black!80, line width=0.8pt, dash pattern=on 4pt off 3pt]
            ([xshift=0.55\textwidth+8pt]frame.north west) --
            ([xshift=0.55\textwidth+8pt]frame.south west);
    }
]
\textbf{Q.18} \quad {\footnotesize\textsc{[NAT $\cdot$ 2 Mark(s) $\cdot$ No Neg]}} \hfill {\footnotesize\textbf{GATE 2023 CS} \quad \tiny\texttt{2023\_CS\_CS\_Q59}}
\vspace{2pt}\hrule height 0.3pt\vspace{4pt}

Consider a sequence $a$ of elements $a_0 = 1, a_1 = 5, a_2 = 7, a_3 = 8, a_4 = 9$, and $a_5 = 2$. The following operations are performed on a stack $S$ and a queue $Q$, both of which are initially empty:
\begin{enumerate}[label=\Roman*.,itemsep=0pt,topsep=2pt]
\item push the elements of $a$ from $a_0$ to $a_5$ in that order into $S$.
\item enqueue the elements of $a$ from $a_0$ to $a_5$ in that order into $Q$.
\item pop an element from $S$.
\item dequeue an element from $Q$.
\item pop an element from $S$.
\item dequeue an element from $Q$.
\item dequeue an element from $Q$ and push the same element into $S$.
\item Repeat operation VII three times.
\item pop an element from $S$.
\item pop an element from $S$.
\end{enumerate}
The top element of $S$ after executing the above operations is \_\_\_\_\_\_\_\_\. (Answer in integer)

\par\vspace{6pt}\noindent\textbf{Your Answer:} \fbox{\rule{0pt}{1.3ex}\hspace{3cm}}

\tcblower

\textbf{Official Key:} \textbf{\large 8 to 8} \hfill \textit{\small [Medium]}
\vspace{2pt}\hrule height 0.3pt\vspace{3pt}
\textbf{Topic:} Programming and Data Structures \\
\textbf{Subtopic:} \textit{Arrays, stacks, queues, linked lists}
\vspace{2pt}\hrule height 0.3pt\vspace{3pt}
\textbf{Guide / Hint:} \\
{\small Operations trace: Push $[1, 5, 7, 8, 9, 2]$ to $S$ (top=2); enqueue $[1, 5, 7, 8, 9, 2]$ to $Q$ (front=1). Pop 2 from $S$ ($S=[1,5,7,8,9]$), dequeue 1 from $Q$ ($Q=[5,7,8,9,2]$); pop 9 from $S$ ($S=[1,5,7,8]$), dequeue 5 from $Q$ ($Q=[7,8,9,2]$). Repeating dequeue from $Q$ and push to $S$ 4 times moves $7, 8, 9, 2$ into $S$, making $S=[1,5,7,8,7,8,9,2]$. Finally, two pops remove 2 and 9, leaving 8 at the top of $S$.}

\end{tcolorbox}
\vspace{6pt}\hrule height 0.35pt\vspace{6pt}
\noindent
\begin{tcolorbox}[
    blanker,
    sidebyside,
    sidebyside align=top,
    lefthand width=0.55\textwidth,
    top=4pt, bottom=4pt, left=0pt, right=0pt,
    middle=16pt,
    underlay={
        \draw[dashed, black!80, line width=0.8pt, dash pattern=on 4pt off 3pt]
            ([xshift=0.55\textwidth+8pt]frame.north west) --
            ([xshift=0.55\textwidth+8pt]frame.south west);
    }
]
\textbf{Q.19} \quad {\footnotesize\textsc{[MCQ $\cdot$ 1 Mark(s) $\cdot$ Neg: 0.33]}} \hfill {\footnotesize\textbf{GATE 2024 CS1} \quad \tiny\texttt{2024\_CS1\_CS\_Q18}}
\vspace{2pt}\hrule height 0.3pt\vspace{4pt}

Consider the following C program:
\begin{lstlisting}[language=C]
#include <stdio.h>
int main() {
    int a = 6;
    int b = 0;
    while(a < 10) {
        a = a / 12 + 1;
        a += b;
    }
    printf("%d", a);
    return 0;
}
\end{lstlisting}
Which one of the following statements is CORRECT?

\par\vspace{6pt}\noindent
\begin{tabular}{@{}p{0.98\linewidth}@{}}
$\square$\ \textbf{(A)} The program prints 9 as output \\
\addlinespace[3pt]
$\square$\ \textbf{(B)} The program prints 10 as output \\
\addlinespace[3pt]
$\square$\ \textbf{(C)} The program gets stuck in an infinite loop \\
\addlinespace[3pt]
$\square$\ \textbf{(D)} The program prints 6 as output
\end{tabular}

\tcblower

\textbf{Official Key:} \textbf{\large C} \hfill \textit{\small [Easy]}
\vspace{2pt}\hrule height 0.3pt\vspace{3pt}
\textbf{Topic:} Programming and Data Structures \\
\textbf{Subtopic:} \textit{Programming in C}
\vspace{2pt}\hrule height 0.3pt\vspace{3pt}
\textbf{Guide / Hint:} \\
{\small Initially $a = 6, b = 0$. In the loop, integer division $a/12 = 6/12 = 0$, so $a = 0 + 1 = 1$, and $a += b \implies a = 1$. In subsequent iterations, $a=1 < 10$, $1/12 + 1 = 1$, so $a$ remains 1 in every iteration. The loop condition $a < 10$ remains true indefinitely, causing an infinite loop.}

\end{tcolorbox}
\vspace{6pt}\hrule height 0.35pt\vspace{6pt}
\noindent
\begin{tcolorbox}[
    blanker,
    sidebyside,
    sidebyside align=top,
    lefthand width=0.55\textwidth,
    top=4pt, bottom=4pt, left=0pt, right=0pt,
    middle=16pt,
    underlay={
        \draw[dashed, black!80, line width=0.8pt, dash pattern=on 4pt off 3pt]
            ([xshift=0.55\textwidth+8pt]frame.north west) --
            ([xshift=0.55\textwidth+8pt]frame.south west);
    }
]
\textbf{Q.20} \quad {\footnotesize\textsc{[MCQ $\cdot$ 1 Mark(s) $\cdot$ Neg: 0.33]}} \hfill {\footnotesize\textbf{GATE 2024 CS1} \quad \tiny\texttt{2024\_CS1\_CS\_Q19}}
\vspace{2pt}\hrule height 0.3pt\vspace{4pt}

Consider the following C program:
\begin{lstlisting}[language=C]
#include <stdio.h>
void fX();
int main() {
    fX();
    return 0;
}
void fX() {
    char a;
    if((a = getchar()) != '\n')
        fX();
    if(a != '\n')
        putchar(a);
}
\end{lstlisting}
Assume that the input to the program from the command line is `1234` followed by a newline character. Which one of the following statements is CORRECT?

\par\vspace{6pt}\noindent
\begin{tabular}{@{}p{0.98\linewidth}@{}}
$\square$\ \textbf{(A)} The program will not terminate \\
\addlinespace[3pt]
$\square$\ \textbf{(B)} The program will terminate with no output \\
\addlinespace[3pt]
$\square$\ \textbf{(C)} The program will terminate with 4321 as output \\
\addlinespace[3pt]
$\square$\ \textbf{(D)} The program will terminate with 1234 as output
\end{tabular}

\tcblower

\textbf{Official Key:} \textbf{\large C} \hfill \textit{\small [Easy]}
\vspace{2pt}\hrule height 0.3pt\vspace{3pt}
\textbf{Topic:} Programming and Data Structures \\
\textbf{Subtopic:} \textit{Recursion}
\vspace{2pt}\hrule height 0.3pt\vspace{3pt}
\textbf{Guide / Hint:} \\
{\small The function `fX()` recursively reads each character before printing. As each character (`'1'`, `'2'`, `'3'`, `'4'`) is read, it is saved on the call stack before `fX()` recurses. When `\n` is reached, recursion terminates and the characters are printed during stack unwinding in reverse (LIFO) order as `4321`.}

\end{tcolorbox}
\vspace{6pt}\hrule height 0.35pt\vspace{6pt}
\noindent
\begin{tcolorbox}[
    blanker,
    sidebyside,
    sidebyside align=top,
    lefthand width=0.55\textwidth,
    top=4pt, bottom=4pt, left=0pt, right=0pt,
    middle=16pt,
    underlay={
        \draw[dashed, black!80, line width=0.8pt, dash pattern=on 4pt off 3pt]
            ([xshift=0.55\textwidth+8pt]frame.north west) --
            ([xshift=0.55\textwidth+8pt]frame.south west);
    }
]
\textbf{Q.21} \quad {\footnotesize\textsc{[MCQ $\cdot$ 2 Mark(s) $\cdot$ Neg: 0.66]}} \hfill {\footnotesize\textbf{GATE 2024 CS1} \quad \tiny\texttt{2024\_CS1\_CS\_Q41}}
\vspace{2pt}\hrule height 0.3pt\vspace{4pt}

An array $[82, 101, 90, 11, 111, 75, 33, 131, 44, 93]$ is heapified. Which one of the following options represents the first three elements in the heapified array?

\par\vspace{6pt}\noindent
\begin{tabular}{@{}p{0.48\linewidth}@{\hspace{0.04\linewidth}}p{0.48\linewidth}@{}}
$\square$\ \textbf{(A)} 82, 90, 101 & $\square$\ \textbf{(B)} 82, 11, 93 \\
\addlinespace[3pt]
$\square$\ \textbf{(C)} 131, 11, 93 & $\square$\ \textbf{(D)} 131, 111, 90
\end{tabular}

\tcblower

\textbf{Official Key:} \textbf{\large D} \hfill \textit{\small [Medium]}
\vspace{2pt}\hrule height 0.3pt\vspace{3pt}
\textbf{Topic:} Programming and Data Structures \\
\textbf{Subtopic:} \textit{Trees, binary search trees, binary heaps}
\vspace{2pt}\hrule height 0.3pt\vspace{3pt}
\textbf{Guide / Hint:} \\
{\small Applying standard bottom-up max-heap construction (Floyd's algorithm) to $[82, 101, 90, 11, 111, 75, 33, 131, 44, 93]$ of length 10: non-leaf nodes are heapified from index 5 down to 1. The maximum element 131 bubbles up to index 1 (root), its left child becomes 111 (index 2), and right child becomes 90 (index 3).}

\end{tcolorbox}
\vspace{6pt}\hrule height 0.35pt\vspace{6pt}
\noindent
\begin{tcolorbox}[
    blanker,
    sidebyside,
    sidebyside align=top,
    lefthand width=0.55\textwidth,
    top=4pt, bottom=4pt, left=0pt, right=0pt,
    middle=16pt,
    underlay={
        \draw[dashed, black!80, line width=0.8pt, dash pattern=on 4pt off 3pt]
            ([xshift=0.55\textwidth+8pt]frame.north west) --
            ([xshift=0.55\textwidth+8pt]frame.south west);
    }
]
\textbf{Q.22} \quad {\footnotesize\textsc{[MCQ $\cdot$ 2 Mark(s) $\cdot$ Neg: 0.66]}} \hfill {\footnotesize\textbf{GATE 2024 CS1} \quad \tiny\texttt{2024\_CS1\_CS\_Q43}}
\vspace{2pt}\hrule height 0.3pt\vspace{4pt}

Consider a binary min-heap containing 105 distinct elements. Let $k$ be the index (in the underlying array) of the maximum element stored in the heap. The number of possible values of $k$ is

\par\vspace{6pt}\noindent
\begin{tabular}{@{}p{0.48\linewidth}@{\hspace{0.04\linewidth}}p{0.48\linewidth}@{}}
$\square$\ \textbf{(A)} 53 & $\square$\ \textbf{(B)} 52 \\
\addlinespace[3pt]
$\square$\ \textbf{(C)} 27 & $\square$\ \textbf{(D)} 1
\end{tabular}

\tcblower

\textbf{Official Key:} \textbf{\large A} \hfill \textit{\small [Medium]}
\vspace{2pt}\hrule height 0.3pt\vspace{3pt}
\textbf{Topic:} Programming and Data Structures \\
\textbf{Subtopic:} \textit{Trees, binary search trees, binary heaps}
\vspace{2pt}\hrule height 0.3pt\vspace{3pt}
\textbf{Guide / Hint:} \\
{\small In a binary min-heap, every parent is strictly smaller than its children, so the maximum element cannot be an internal node and must reside at one of the leaf nodes. In a complete binary tree of $n = 105$ nodes, the number of leaf nodes is $\lceil n/2 \rceil = \lceil 105/2 \rceil = 53$. Hence, there are 53 possible index locations for the maximum element.}

\end{tcolorbox}
\vspace{6pt}\hrule height 0.35pt\vspace{6pt}
\noindent
\begin{tcolorbox}[
    blanker,
    sidebyside,
    sidebyside align=top,
    lefthand width=0.55\textwidth,
    top=4pt, bottom=4pt, left=0pt, right=0pt,
    middle=16pt,
    underlay={
        \draw[dashed, black!80, line width=0.8pt, dash pattern=on 4pt off 3pt]
            ([xshift=0.55\textwidth+8pt]frame.north west) --
            ([xshift=0.55\textwidth+8pt]frame.south west);
    }
]
\textbf{Q.23} \quad {\footnotesize\textsc{[MSQ $\cdot$ 2 Mark(s) $\cdot$ No Neg]}} \hfill {\footnotesize\textbf{GATE 2024 CS1} \quad \tiny\texttt{2024\_CS1\_CS\_Q48}}
\vspace{2pt}\hrule height 0.3pt\vspace{4pt}

Consider the following C function definition:
\begin{lstlisting}[language=C]
int f(int x, int y) {
    for (int i = 0; i < y; i++) {
        x = x + x + y;
    }
    return x;
}
\end{lstlisting}
Which of the following statements is/are TRUE about the above function?

\par\vspace{6pt}\noindent
\begin{tabular}{@{}p{0.98\linewidth}@{}}
$\square$\ \textbf{(A)} If the inputs are $x=20, y=10$, then the return value is greater than $2^{20}$ \\
\addlinespace[3pt]
$\square$\ \textbf{(B)} If the inputs are $x=20, y=20$, then the return value is greater than $2^{20}$ \\
\addlinespace[3pt]
$\square$\ \textbf{(C)} If the inputs are $x=20, y=10$, then the return value is less than $2^{10}$ \\
\addlinespace[3pt]
$\square$\ \textbf{(D)} If the inputs are $x=10, y=20$, then the return value is greater than $2^{20}$
\end{tabular}

\tcblower

\textbf{Official Key:} \textbf{\large B;D} \hfill \textit{\small [Medium]}
\vspace{2pt}\hrule height 0.3pt\vspace{3pt}
\textbf{Topic:} Programming and Data Structures \\
\textbf{Subtopic:} \textit{Programming in C}
\vspace{2pt}\hrule height 0.3pt\vspace{3pt}
\textbf{Guide / Hint:} \\
{\small In each loop iteration, $x_{k+1} = 2x_k + y$. Solving this linear recurrence yields $x_y = 2^y x_0 + y(2^y - 1) = 2^y(x_0 + y) - y$. For $x=20, y=20$: $x_{20} = 2^{20}(40) - 20 > 2^{20}$ (B is TRUE). For $x=10, y=20$: $x_{20} = 2^{20}(30) - 20 > 2^{20}$ (D is TRUE). For $y=10, x=20$: $x_{10} = 2^{10}(30) - 10 = 30710 > 2^{10}$, so (A) and (C) are false.}

\end{tcolorbox}
\vspace{6pt}\hrule height 0.35pt\vspace{6pt}
\noindent
\begin{tcolorbox}[
    blanker,
    sidebyside,
    sidebyside align=top,
    lefthand width=0.55\textwidth,
    top=4pt, bottom=4pt, left=0pt, right=0pt,
    middle=16pt,
    underlay={
        \draw[dashed, black!80, line width=0.8pt, dash pattern=on 4pt off 3pt]
            ([xshift=0.55\textwidth+8pt]frame.north west) --
            ([xshift=0.55\textwidth+8pt]frame.south west);
    }
]
\textbf{Q.24} \quad {\footnotesize\textsc{[MCQ $\cdot$ 1 Mark(s) $\cdot$ Neg: 0.33]}} \hfill {\footnotesize\textbf{GATE 2024 CS2} \quad \tiny\texttt{2024\_CS2\_CS\_Q13}}
\vspace{2pt}\hrule height 0.3pt\vspace{4pt}

Consider the following C program. Assume parameters to a function are evaluated from right to left.
\begin{verbatim}
#include <stdio.h>
int g(int p) { printf("%d", p); return p; }
int h(int q) { printf("%d", q); return q; }
void f(int x, int y) {
    g(x);
    h(y);
}
int main() {
    f(g(10), h(20));
}
\end{verbatim}
Which one of the following is the output of the program?

\par\vspace{6pt}\noindent
\begin{tabular}{@{}p{0.48\linewidth}@{\hspace{0.04\linewidth}}p{0.48\linewidth}@{}}
$\square$\ \textbf{(A)} 20101020 & $\square$\ \textbf{(B)} 10202010 \\
\addlinespace[3pt]
$\square$\ \textbf{(C)} 20102010 & $\square$\ \textbf{(D)} 10201020
\end{tabular}

\tcblower

\textbf{Official Key:} \textbf{\large A} \hfill \textit{\small [Medium]}
\vspace{2pt}\hrule height 0.3pt\vspace{3pt}
\textbf{Topic:} Programming and Data Structures \\
\textbf{Subtopic:} \textit{Programming in C}
\vspace{2pt}\hrule height 0.3pt\vspace{3pt}
\textbf{Guide / Hint:} \\
{\small In `main()`, `f(g(10), h(20))` evaluates parameters from right to left: `h(20)` executes first, printing `20` and returning 20 ($y = 20$), then `g(10)` executes, printing `10` and returning 10 ($x = 10$). Inside `f`, `g(x)` executes printing `10`, followed by `h(y)` printing `20`. The complete printed sequence is `20101020`.}

\end{tcolorbox}
\vspace{6pt}\hrule height 0.35pt\vspace{6pt}
\noindent
\begin{tcolorbox}[
    blanker,
    sidebyside,
    sidebyside align=top,
    lefthand width=0.55\textwidth,
    top=4pt, bottom=4pt, left=0pt, right=0pt,
    middle=16pt,
    underlay={
        \draw[dashed, black!80, line width=0.8pt, dash pattern=on 4pt off 3pt]
            ([xshift=0.55\textwidth+8pt]frame.north west) --
            ([xshift=0.55\textwidth+8pt]frame.south west);
    }
]
\textbf{Q.25} \quad {\footnotesize\textsc{[MSQ $\cdot$ 1 Mark(s) $\cdot$ No Neg]}} \hfill {\footnotesize\textbf{GATE 2024 CS2} \quad \tiny\texttt{2024\_CS2\_CS\_Q33}}
\vspace{2pt}\hrule height 0.3pt\vspace{4pt}

Consider the following C function definition:
\begin{verbatim}
int fX(char *a){
    char *b = a;
    while(*b)
        b++;
    return b - a;
}
\end{verbatim}
Which of the following statements is/are TRUE?

\par\vspace{6pt}\noindent
\begin{tabular}{@{}p{0.98\linewidth}@{}}
$\square$\ \textbf{(A)} The function call fX("abcd") will always return a value \\
\addlinespace[3pt]
$\square$\ \textbf{(B)} Assuming a character array c is declared as char c[] = "abcd" in main(), the function call fX(c) will always return a value \\
\addlinespace[3pt]
$\square$\ \textbf{(C)} The code of the function will not compile \\
\addlinespace[3pt]
$\square$\ \textbf{(D)} Assuming a character pointer c is declared as char *c = "abcd" in main(), the function call fX(c) will always return a value
\end{tabular}

\tcblower

\textbf{Official Key:} \textbf{\large A;B;D} \hfill \textit{\small [Easy]}
\vspace{2pt}\hrule height 0.3pt\vspace{3pt}
\textbf{Topic:} Programming and Data Structures \\
\textbf{Subtopic:} \textit{Programming in C}
\vspace{2pt}\hrule height 0.3pt\vspace{3pt}
\textbf{Guide / Hint:} \\
{\small The function `fX` calculates the string length by advancing pointer `b` until the null terminator `\0` is reached, then returning pointer difference `b - a`. String literals in C are null-terminated, so calls with string literals or character arrays (`fX("abcd")`, `fX(c)`) correctly compute 4 and return (A, B, and D are true).}

\end{tcolorbox}
\vspace{6pt}\hrule height 0.35pt\vspace{6pt}
\noindent
\begin{tcolorbox}[
    blanker,
    sidebyside,
    sidebyside align=top,
    lefthand width=0.55\textwidth,
    top=4pt, bottom=4pt, left=0pt, right=0pt,
    middle=16pt,
    underlay={
        \draw[dashed, black!80, line width=0.8pt, dash pattern=on 4pt off 3pt]
            ([xshift=0.55\textwidth+8pt]frame.north west) --
            ([xshift=0.55\textwidth+8pt]frame.south west);
    }
]
\textbf{Q.26} \quad {\footnotesize\textsc{[MCQ $\cdot$ 2 Mark(s) $\cdot$ Neg: 0.66]}} \hfill {\footnotesize\textbf{GATE 2024 CS2} \quad \tiny\texttt{2024\_CS2\_CS\_Q36}}
\vspace{2pt}\hrule height 0.3pt\vspace{4pt}

What is the output of the following C program?
\begin{verbatim}
#include <stdio.h>
int main() {
    double a[2] = {20.0, 25.0}, *p, *q;
    p = a;
    q = p + 1;
    printf("%d,%d", (int)(q - p), (int)(*q - *p));
    return 0;
}
\end{verbatim}

\par\vspace{6pt}\noindent
\begin{tabular}{@{}p{0.48\linewidth}@{\hspace{0.04\linewidth}}p{0.48\linewidth}@{}}
$\square$\ \textbf{(A)} 4,8 & $\square$\ \textbf{(B)} 1,5 \\
\addlinespace[3pt]
$\square$\ \textbf{(C)} 8,5 & $\square$\ \textbf{(D)} 1,8
\end{tabular}

\tcblower

\textbf{Official Key:} \textbf{\large B} \hfill \textit{\small [Easy]}
\vspace{2pt}\hrule height 0.3pt\vspace{3pt}
\textbf{Topic:} Programming and Data Structures \\
\textbf{Subtopic:} \textit{Programming in C}
\vspace{2pt}\hrule height 0.3pt\vspace{3pt}
\textbf{Guide / Hint:} \\
{\small In pointer arithmetic, pointer subtraction `q - p` yields the number of elements between pointers, so $(q - p) = 1$. The expression `*q - *p` calculates the difference between pointed double values: $25.0 - 20.0 = 5.0$, cast to integer is 5. Thus, `1,5` is printed.}

\end{tcolorbox}
\vspace{6pt}\hrule height 0.35pt\vspace{6pt}
\noindent
\begin{tcolorbox}[
    blanker,
    sidebyside,
    sidebyside align=top,
    lefthand width=0.55\textwidth,
    top=4pt, bottom=4pt, left=0pt, right=0pt,
    middle=16pt,
    underlay={
        \draw[dashed, black!80, line width=0.8pt, dash pattern=on 4pt off 3pt]
            ([xshift=0.55\textwidth+8pt]frame.north west) --
            ([xshift=0.55\textwidth+8pt]frame.south west);
    }
]
\textbf{Q.27} \quad {\footnotesize\textsc{[MCQ $\cdot$ 2 Mark(s) $\cdot$ Neg: 0.66]}} \hfill {\footnotesize\textbf{GATE 2024 CS2} \quad \tiny\texttt{2024\_CS2\_CS\_Q39}}
\vspace{2pt}\hrule height 0.3pt\vspace{4pt}

You are given a set $V$ of distinct integers. A binary search tree $T$ is created by inserting all elements of $V$ one by one, starting with an empty tree. The tree $T$ follows the convention that, at each node, all values stored in the left subtree of the node are smaller than the value stored at the node. Which of the following is TRUE?

\par\vspace{6pt}\noindent
\begin{tabular}{@{}p{0.98\linewidth}@{}}
$\square$\ \textbf{(A)} Inorder traversal of $T$ can be determined from $V$ \\
\addlinespace[3pt]
$\square$\ \textbf{(B)} Root node of $T$ can be determined from $V$ \\
\addlinespace[3pt]
$\square$\ \textbf{(C)} Preorder traversal of $T$ can be determined from $V$ \\
\addlinespace[3pt]
$\square$\ \textbf{(D)} Postorder traversal of $T$ can be determined from $V$
\end{tabular}

\tcblower

\textbf{Official Key:} \textbf{\large A} \hfill \textit{\small [Easy]}
\vspace{2pt}\hrule height 0.3pt\vspace{3pt}
\textbf{Topic:} Programming and Data Structures \\
\textbf{Subtopic:} \textit{Trees, binary search trees, binary heaps}
\vspace{2pt}\hrule height 0.3pt\vspace{3pt}
\textbf{Guide / Hint:} \\
{\small For any binary search tree, the inorder traversal sequence always visits elements in strictly increasing sorted order regardless of insertion order or tree shape. Since the set of distinct elements $V$ is known, sorting $V$ directly determines the unique inorder traversal.}

\end{tcolorbox}
\vspace{6pt}\hrule height 0.35pt\vspace{6pt}
\noindent
\begin{tcolorbox}[
    blanker,
    sidebyside,
    sidebyside align=top,
    lefthand width=0.55\textwidth,
    top=4pt, bottom=4pt, left=0pt, right=0pt,
    middle=16pt,
    underlay={
        \draw[dashed, black!80, line width=0.8pt, dash pattern=on 4pt off 3pt]
            ([xshift=0.55\textwidth+8pt]frame.north west) --
            ([xshift=0.55\textwidth+8pt]frame.south west);
    }
]
\textbf{Q.28} \quad {\footnotesize\textsc{[MCQ $\cdot$ 2 Mark(s) $\cdot$ Neg: 0.66]}} \hfill {\footnotesize\textbf{GATE 2024 CS2} \quad \tiny\texttt{2024\_CS2\_CS\_Q42}}
\vspace{2pt}\hrule height 0.3pt\vspace{4pt}

Consider an array X that contains n positive integers. A subarray of X is defined to be a sequence of array locations with consecutive indices.
The C code snippet given below has been written to compute the length of the longest subarray of X that contains at most two distinct integers. The code has two missing expressions labelled $(P)$ and $(Q)$:
\begin{verbatim}
int first=0, second=0, len1=0, len2=0, maxlen=0;
for (int i=0; i < n; i++) {
    if (X[i] == first) {
        len2++; len1++;
    } else if (X[i] == second) {
        len2++;
        len1 = (P);
        second = first;
    } else {
        len2 = (Q);
        len1 = 1; second = first;
    }
    if (len2 > maxlen) maxlen = len2;
    first = X[i];
}
\end{verbatim}
Which one of the following options gives the CORRECT missing expressions?

\par\vspace{6pt}\noindent
\begin{tabular}{@{}p{0.48\linewidth}@{\hspace{0.04\linewidth}}p{0.48\linewidth}@{}}
$\square$\ \textbf{(A)} (P) len1+1 (Q) len2+1 & $\square$\ \textbf{(B)} (P) 1 (Q) len1+1 \\
\addlinespace[3pt]
$\square$\ \textbf{(C)} (P) 1 (Q) len2+1 & $\square$\ \textbf{(D)} (P) len2+1 (Q) len1+1
\end{tabular}

\tcblower

\textbf{Official Key:} \textbf{\large B} \hfill \textit{\small [Hard]}
\vspace{2pt}\hrule height 0.3pt\vspace{3pt}
\textbf{Topic:} Programming and Data Structures \\
\textbf{Subtopic:} \textit{Programming in C}
\vspace{2pt}\hrule height 0.3pt\vspace{3pt}
\textbf{Guide / Hint:} \\
{\small When $X[i] == second$, the current element matches the older distinct value, switching identity with `first`. Because $X[i]$ has just appeared once consecutively, `len1` resets to 1 (expression $P = 1$). When encountering a third distinct value (else branch), the previous window keeping at most two distinct values must keep only the contiguous run of `first` plus current element, so $Q = len1 + 1$.}

\end{tcolorbox}
\vspace{6pt}\hrule height 0.35pt\vspace{6pt}
\noindent
\begin{tcolorbox}[
    blanker,
    sidebyside,
    sidebyside align=top,
    lefthand width=0.55\textwidth,
    top=4pt, bottom=4pt, left=0pt, right=0pt,
    middle=16pt,
    underlay={
        \draw[dashed, black!80, line width=0.8pt, dash pattern=on 4pt off 3pt]
            ([xshift=0.55\textwidth+8pt]frame.north west) --
            ([xshift=0.55\textwidth+8pt]frame.south west);
    }
]
\textbf{Q.29} \quad {\footnotesize\textsc{[MSQ $\cdot$ 2 Mark(s) $\cdot$ No Neg]}} \hfill {\footnotesize\textbf{GATE 2024 CS2} \quad \tiny\texttt{2024\_CS2\_CS\_Q48}}
\vspace{2pt}\hrule height 0.3pt\vspace{4pt}

Let S1 and S2 be two stacks. S1 has capacity of 4 elements. S2 has capacity of 2 elements. S1 already has 4 elements: 100, 200, 300, and 400 (top), whereas S2 is empty.
Only the following three operations are available: PushToS2, PushToS1, and GenerateOutput.
Which of the following output sequences can be generated?

\par\vspace{6pt}\noindent
\begin{tabular}{@{}p{0.48\linewidth}@{\hspace{0.04\linewidth}}p{0.48\linewidth}@{}}
$\square$\ \textbf{(A)} 100, 200, 400, 300 & $\square$\ \textbf{(B)} 200, 300, 400, 100 \\
\addlinespace[3pt]
$\square$\ \textbf{(C)} 400, 200, 100, 300 & $\square$\ \textbf{(D)} 300, 200, 400, 100
\end{tabular}

\tcblower

\textbf{Official Key:} \textbf{\large B;C;D} \hfill \textit{\small [Medium]}
\vspace{2pt}\hrule height 0.3pt\vspace{3pt}
\textbf{Topic:} Programming and Data Structures \\
\textbf{Subtopic:} \textit{Arrays, stacks, queues, linked lists}
\vspace{2pt}\hrule height 0.3pt\vspace{3pt}
\textbf{Guide / Hint:} \\
{\small S1 initially contains [100, 200, 300, 400] from bottom to top; S2 has capacity 2. Sequence (B): Push 400 to S2, push 300 to S2, output 200, pop 300 to S1 and output 300, pop 400 to S1 and output 400, output 100. Similarly, sequences (C) and (D) can be generated within S2's 2-element capacity limit.}

\end{tcolorbox}
\vspace{6pt}\hrule height 0.35pt\vspace{6pt}
\noindent
\begin{tcolorbox}[
    blanker,
    sidebyside,
    sidebyside align=top,
    lefthand width=0.55\textwidth,
    top=4pt, bottom=4pt, left=0pt, right=0pt,
    middle=16pt,
    underlay={
        \draw[dashed, black!80, line width=0.8pt, dash pattern=on 4pt off 3pt]
            ([xshift=0.55\textwidth+8pt]frame.north west) --
            ([xshift=0.55\textwidth+8pt]frame.south west);
    }
]
\textbf{Q.30} \quad {\footnotesize\textsc{[MSQ $\cdot$ 1 Mark(s) $\cdot$ No Neg]}} \hfill {\footnotesize\textbf{GATE 2025 CS1} \quad \tiny\texttt{2025\_CS1\_CS\_Q26}}
\vspace{2pt}\hrule height 0.3pt\vspace{4pt}

Which of the following statement(s) is/are TRUE for any binary search tree (BST) having $n$ distinct integers?

\par\vspace{6pt}\noindent
\begin{tabular}{@{}p{0.98\linewidth}@{}}
$\square$\ \textbf{(A)} The maximum length of a path from the root node to any other node is $(n-1)$. \\
\addlinespace[3pt]
$\square$\ \textbf{(B)} An inorder traversal will always produce a sorted sequence of elements. \\
\addlinespace[3pt]
$\square$\ \textbf{(C)} Finding an element takes $O(\log_2 n)$ time in the worst case. \\
\addlinespace[3pt]
$\square$\ \textbf{(D)} Every BST is also a Min-Heap.
\end{tabular}

\tcblower

\textbf{Official Key:} \textbf{\large A;B} \hfill \textit{\small [Easy]}
\vspace{2pt}\hrule height 0.3pt\vspace{3pt}
\textbf{Topic:} Programming and Data Structures \\
\textbf{Subtopic:} \textit{Trees, binary search trees, binary heaps}
\vspace{2pt}\hrule height 0.3pt\vspace{3pt}
\textbf{Guide / Hint:} \\
{\small In a degenerate (skewed) binary search tree, the path length from root to the farthest leaf is $n - 1$ edges (A is true). Inorder traversal of any valid BST visits nodes in strictly increasing order (B is true). Worst-case search in an unbalanced BST takes $O(n)$ time, and BST property differs fundamentally from heap order.}

\end{tcolorbox}
\vspace{6pt}\hrule height 0.35pt\vspace{6pt}
\noindent
\begin{tcolorbox}[
    blanker,
    sidebyside,
    sidebyside align=top,
    lefthand width=0.55\textwidth,
    top=4pt, bottom=4pt, left=0pt, right=0pt,
    middle=16pt,
    underlay={
        \draw[dashed, black!80, line width=0.8pt, dash pattern=on 4pt off 3pt]
            ([xshift=0.55\textwidth+8pt]frame.north west) --
            ([xshift=0.55\textwidth+8pt]frame.south west);
    }
]
\textbf{Q.31} \quad {\footnotesize\textsc{[NAT $\cdot$ 1 Mark(s) $\cdot$ No Neg]}} \hfill {\footnotesize\textbf{GATE 2025 CS1} \quad \tiny\texttt{2025\_CS1\_CS\_Q34}}
\vspace{2pt}\hrule height 0.3pt\vspace{4pt}

Consider the following C program:
\par\medskip
\begin{verbatim}
#include <stdio.h>
void foo(int *p, int x){
    *p = x;
}
int main(){
    int *z;
    int a = 20, b = 25;
    z = &a;
    foo(z, b);
    printf("%d", a);
    return 0;
}
\end{verbatim}
The output of the given C program is \_\_\_\_\_\_\. (Answer in integer)

\par\vspace{6pt}\noindent\textbf{Your Answer:} \fbox{\rule{0pt}{1.3ex}\hspace{3cm}}

\tcblower

\textbf{Official Key:} \textbf{\large 25 to 25} \hfill \textit{\small [Easy]}
\vspace{2pt}\hrule height 0.3pt\vspace{3pt}
\textbf{Topic:} Programming and Data Structures \\
\textbf{Subtopic:} \textit{Programming in C}
\vspace{2pt}\hrule height 0.3pt\vspace{3pt}
\textbf{Guide / Hint:} \\
{\small In `main()`, pointer `z` is assigned the memory address of variable `a` (`\&a`). Function `foo(z, b)` dereferences `p` (which points to `a`) and assigns `x = b = 25` to `*p`. Thus, `a` is updated to 25, which `printf` prints.}

\end{tcolorbox}
\vspace{6pt}\hrule height 0.35pt\vspace{6pt}
\noindent
\begin{tcolorbox}[
    blanker,
    sidebyside,
    sidebyside align=top,
    lefthand width=0.55\textwidth,
    top=4pt, bottom=4pt, left=0pt, right=0pt,
    middle=16pt,
    underlay={
        \draw[dashed, black!80, line width=0.8pt, dash pattern=on 4pt off 3pt]
            ([xshift=0.55\textwidth+8pt]frame.north west) --
            ([xshift=0.55\textwidth+8pt]frame.south west);
    }
]
\textbf{Q.32} \quad {\footnotesize\textsc{[NAT $\cdot$ 1 Mark(s) $\cdot$ No Neg]}} \hfill {\footnotesize\textbf{GATE 2025 CS1} \quad \tiny\texttt{2025\_CS1\_CS\_Q35}}
\vspace{2pt}\hrule height 0.3pt\vspace{4pt}

The height of any rooted tree is defined as the maximum number of edges in the path from the root node to any leaf node. Suppose a Min-Heap $T$ stores 32 keys. The height of $T$ is \_\_\_\_\_\_\. (Answer in integer)

\par\vspace{6pt}\noindent\textbf{Your Answer:} \fbox{\rule{0pt}{1.3ex}\hspace{3cm}}

\tcblower

\textbf{Official Key:} \textbf{\large 5 to 5} \hfill \textit{\small [Easy]}
\vspace{2pt}\hrule height 0.3pt\vspace{3pt}
\textbf{Topic:} Programming and Data Structures \\
\textbf{Subtopic:} \textit{Trees, binary search trees, binary heaps}
\vspace{2pt}\hrule height 0.3pt\vspace{3pt}
\textbf{Guide / Hint:} \\
{\small A binary heap storing $n$ keys is an almost complete binary tree whose height in terms of edges is $\lfloor \log_2 n \rfloor$. For $n = 32$, the height is $\lfloor \log_2 32 \rfloor = 5$.}

\end{tcolorbox}
\vspace{6pt}\hrule height 0.35pt\vspace{6pt}
\noindent
\begin{tcolorbox}[
    blanker,
    sidebyside,
    sidebyside align=top,
    lefthand width=0.55\textwidth,
    top=4pt, bottom=4pt, left=0pt, right=0pt,
    middle=16pt,
    underlay={
        \draw[dashed, black!80, line width=0.8pt, dash pattern=on 4pt off 3pt]
            ([xshift=0.55\textwidth+8pt]frame.north west) --
            ([xshift=0.55\textwidth+8pt]frame.south west);
    }
]
\textbf{Q.33} \quad {\footnotesize\textsc{[NAT $\cdot$ 2 Mark(s) $\cdot$ No Neg]}} \hfill {\footnotesize\textbf{GATE 2025 CS1} \quad \tiny\texttt{2025\_CS1\_CS\_Q61}}
\vspace{2pt}\hrule height 0.3pt\vspace{4pt}

Consider the following C program:
\par\medskip
\begin{verbatim}
#include <stdio.h>
int foo(int S[], int size){
    if(size == 0) return 0;
    if(size == 1) return 1;
    if(S[0] != S[1]) return 1 + foo(S+1, size-1);
    return foo(S+1, size-1);
}
int main(){
    int A[] = {0, 1, 2, 2, 2, 0, 0, 1, 1};
    printf("%d", foo(A, 9));
    return 0;
}
\end{verbatim}
The value printed by the given C program is \_\_\_\_\_\_\. (Answer in integer)

\par\vspace{6pt}\noindent\textbf{Your Answer:} \fbox{\rule{0pt}{1.3ex}\hspace{3cm}}

\tcblower

\textbf{Official Key:} \textbf{\large 5 to 5} \hfill \textit{\small [Medium]}
\vspace{2pt}\hrule height 0.3pt\vspace{3pt}
\textbf{Topic:} Programming and Data Structures \\
\textbf{Subtopic:} \textit{Recursion}
\vspace{2pt}\hrule height 0.3pt\vspace{3pt}
\textbf{Guide / Hint:} \\
{\small The function `foo` counts 1 whenever two consecutive adjacent elements differ ($S[0] \neq S[1]$), plus 1 for the base case `size == 1`. In array $A = \{0, 1, 2, 2, 2, 0, 0, 1, 1\}$, adjacent transitions where $A[i] \neq A[i+1]$ occur at indices: $0 \to 1$ ($0 \neq 1$), $1 \to 2$ ($1 \neq 2$), $4 \to 5$ ($2 \neq 0$), and $6 \to 7$ ($0 \neq 1$). There are 4 such differences, plus 1 returned at `size == 1`, giving $4 + 1 = 5$.}

\end{tcolorbox}
\vspace{6pt}\hrule height 0.35pt\vspace{6pt}
\noindent
\begin{tcolorbox}[
    blanker,
    sidebyside,
    sidebyside align=top,
    lefthand width=0.55\textwidth,
    top=4pt, bottom=4pt, left=0pt, right=0pt,
    middle=16pt,
    underlay={
        \draw[dashed, black!80, line width=0.8pt, dash pattern=on 4pt off 3pt]
            ([xshift=0.55\textwidth+8pt]frame.north west) --
            ([xshift=0.55\textwidth+8pt]frame.south west);
    }
]
\textbf{Q.34} \quad {\footnotesize\textsc{[NAT $\cdot$ 2 Mark(s) $\cdot$ No Neg]}} \hfill {\footnotesize\textbf{GATE 2025 CS1} \quad \tiny\texttt{2025\_CS1\_CS\_Q62}}
\vspace{2pt}\hrule height 0.3pt\vspace{4pt}

Let LIST be a datatype for an implementation of linked list defined as follows:
\begin{verbatim}
typedef struct list {
    int data;
    struct list *next;
} LIST;
\end{verbatim}
Suppose a program has created two linked lists, L1 and L2, whose contents are given (code for creating L1 and L2 is not provided here). L1 contains 9 nodes, and L2 contains 7 nodes.
Consider the following C program segment that modifies the list L1. The number of nodes that will be there in L1 after the execution of the code segment is \_\_\_\_\_\_\. (Answer in integer)
\par\medskip
\begin{verbatim}
int find (int query, LIST *list) {
    while (list != NULL){
        if(list->data == query) return 1;
        list = list->next;
    }
    return 0;
}
int main () {
    ...
    ptr1 = L1; ptr2 = L2;
    while (ptr1->next != NULL){
        query = ptr1->next->data;
        if (find (query, L2))
            ptr1->next = ptr1->next->next;
        else ptr1 = ptr1->next;
    }
    ...
    return 0;
}
\end{verbatim}

\par\vspace{6pt}\noindent\textbf{Your Answer:} \fbox{\rule{0pt}{1.3ex}\hspace{3cm}}

\tcblower

\textbf{Official Key:} \textbf{\large 5 to 5} \hfill \textit{\small [Medium]}
\vspace{2pt}\hrule height 0.3pt\vspace{3pt}
\textbf{Topic:} Programming and Data Structures \\
\textbf{Subtopic:} \textit{Arrays, stacks, queues, linked lists}
\vspace{2pt}\hrule height 0.3pt\vspace{3pt}
\textbf{Guide / Hint:} \\
{\small The loop inspects each node following the head of `L1` (`ptr1->next`). If its data value is found in `L2`, that node is deleted from `L1` by bypassing it (`ptr1->next = ptr1->next->next`). Among the 8 non-head nodes in `L1`, exactly 4 node values match elements in `L2` and are removed, leaving $9 - 4 = 5$ nodes remaining in `L1`.}

\end{tcolorbox}
\vspace{6pt}\hrule height 0.35pt\vspace{6pt}
\noindent
\begin{tcolorbox}[
    blanker,
    sidebyside,
    sidebyside align=top,
    lefthand width=0.55\textwidth,
    top=4pt, bottom=4pt, left=0pt, right=0pt,
    middle=16pt,
    underlay={
        \draw[dashed, black!80, line width=0.8pt, dash pattern=on 4pt off 3pt]
            ([xshift=0.55\textwidth+8pt]frame.north west) --
            ([xshift=0.55\textwidth+8pt]frame.south west);
    }
]
\textbf{Q.35} \quad {\footnotesize\textsc{[NAT $\cdot$ 2 Mark(s) $\cdot$ No Neg]}} \hfill {\footnotesize\textbf{GATE 2025 CS1} \quad \tiny\texttt{2025\_CS1\_CS\_Q63}}
\vspace{2pt}\hrule height 0.3pt\vspace{4pt}

Consider the following C program:
\par\medskip
\begin{verbatim}
#include <stdio.h>
int gate (int n) {
    int d, t, newnum, turn;
    newnum = turn = 0; t = 1;
    while (n >= t) t *= 10;
    t /= 10;
    while (t > 0) {
        d = n / t;
        n = n % t;
        t /= 10;
        if (turn) newnum = 10 * newnum + d;
        turn = (turn + 1) % 2;
    }
    return newnum;
}
int main () {
    printf("%d", gate(14362));
    return 0;
}
\end{verbatim}
The value printed by the given C program is \_\_\_\_\_\_\. (Answer in integer)

\par\vspace{6pt}\noindent\textbf{Your Answer:} \fbox{\rule{0pt}{1.3ex}\hspace{3cm}}

\tcblower

\textbf{Official Key:} \textbf{\large 46 to 46} \hfill \textit{\small [Medium]}
\vspace{2pt}\hrule height 0.3pt\vspace{3pt}
\textbf{Topic:} Programming and Data Structures \\
\textbf{Subtopic:} \textit{Programming in C}
\vspace{2pt}\hrule height 0.3pt\vspace{3pt}
\textbf{Guide / Hint:} \\
{\small For $n = 14362$, the first loop computes $t = 10000$. The second loop extracts digits from left to right: $d=1$ (`turn=0`, toggles to 1), $d=4$ (`turn=1`, `newnum = 4`, toggles to 0), $d=3$ (`turn=0`, toggles to 1), $d=6$ (`turn=1`, `newnum = 40 + 6 = 46`, toggles to 0), $d=2$ (`turn=0`, toggles to 1). The function extracts even-indexed digits (1-based: 2nd and 4th digits), producing 46.}

\end{tcolorbox}
\vspace{6pt}\hrule height 0.35pt\vspace{6pt}
\noindent
\begin{tcolorbox}[
    blanker,
    sidebyside,
    sidebyside align=top,
    lefthand width=0.55\textwidth,
    top=4pt, bottom=4pt, left=0pt, right=0pt,
    middle=16pt,
    underlay={
        \draw[dashed, black!80, line width=0.8pt, dash pattern=on 4pt off 3pt]
            ([xshift=0.55\textwidth+8pt]frame.north west) --
            ([xshift=0.55\textwidth+8pt]frame.south west);
    }
]
\textbf{Q.36} \quad {\footnotesize\textsc{[MCQ $\cdot$ 1 Mark(s) $\cdot$ Neg: 0.33]}} \hfill {\footnotesize\textbf{GATE 2025 CS2} \quad \tiny\texttt{2025\_CS2\_CS\_Q13}}
\vspace{2pt}\hrule height 0.3pt\vspace{4pt}

Consider a binary tree $T$ in which every node has either zero or two children. Let $n > 0$ be the number of nodes in $T$. Which ONE of the following is the number of nodes in $T$ that have exactly two children?

\par\vspace{6pt}\noindent
\begin{tabular}{@{}p{0.48\linewidth}@{\hspace{0.04\linewidth}}p{0.48\linewidth}@{}}
$\square$\ \textbf{(A)} $\frac{n-2}{2}$ & $\square$\ \textbf{(B)} $\frac{n-1}{2}$ \\
\addlinespace[3pt]
$\square$\ \textbf{(C)} $\frac{n}{2}$ & $\square$\ \textbf{(D)} $\frac{n+1}{2}$
\end{tabular}

\tcblower

\textbf{Official Key:} \textbf{\large B} \hfill \textit{\small [Easy]}
\vspace{2pt}\hrule height 0.3pt\vspace{3pt}
\textbf{Topic:} Programming and Data Structures \\
\textbf{Subtopic:} \textit{Trees and strictly binary trees}
\vspace{2pt}\hrule height 0.3pt\vspace{3pt}
\textbf{Guide / Hint:} \\
{\small In any proper (strictly) binary tree with $L$ leaves and $I$ internal nodes (nodes with 2 children), $L = I + 1$. Since $n = L + I = 2I + 1$, the number of internal nodes is $I = \frac{n-1}{2}$.}

\end{tcolorbox}
\vspace{6pt}\hrule height 0.35pt\vspace{6pt}
\noindent
\begin{tcolorbox}[
    blanker,
    sidebyside,
    sidebyside align=top,
    lefthand width=0.55\textwidth,
    top=4pt, bottom=4pt, left=0pt, right=0pt,
    middle=16pt,
    underlay={
        \draw[dashed, black!80, line width=0.8pt, dash pattern=on 4pt off 3pt]
            ([xshift=0.55\textwidth+8pt]frame.north west) --
            ([xshift=0.55\textwidth+8pt]frame.south west);
    }
]
\textbf{Q.37} \quad {\footnotesize\textsc{[MCQ $\cdot$ 1 Mark(s) $\cdot$ Neg: 0.33]}} \hfill {\footnotesize\textbf{GATE 2025 CS2} \quad \tiny\texttt{2025\_CS2\_CS\_Q19}}
\vspace{2pt}\hrule height 0.3pt\vspace{4pt}

Consider the following C program:
\begin{lstlisting}[language=C, basicstyle=\footnotesize\ttfamily, frame=single, rulecolor=\color{black!30}, tabsize=2, breaklines=true]
#include <stdio.h>
void stringcopy(char *, char *);
int main() {
    char a[30] = "@#Hello World!";
    stringcopy(a, a + 2);
    printf("%s\n", a);
    return 0;
}
void stringcopy(char *s, char *t) {
    while(*t)
        *s++ = *t++;
}
\end{lstlisting}
Which ONE of the following will be the output of the program?

\par\vspace{6pt}\noindent
\begin{tabular}{@{}p{0.48\linewidth}@{\hspace{0.04\linewidth}}p{0.48\linewidth}@{}}
$\square$\ \textbf{(A)} @\#Hello World! & $\square$\ \textbf{(B)} Hello World! \\
\addlinespace[3pt]
$\square$\ \textbf{(C)} ello World! & $\square$\ \textbf{(D)} Hello World!d!
\end{tabular}

\tcblower

\textbf{Official Key:} \textbf{\large D} \hfill \textit{\small [Medium]}
\vspace{2pt}\hrule height 0.3pt\vspace{3pt}
\textbf{Topic:} Programming and Data Structures \\
\textbf{Subtopic:} \textit{Programming in C: pointers and strings}
\vspace{2pt}\hrule height 0.3pt\vspace{3pt}
\textbf{Guide / Hint:} \\
{\small `stringcopy(a, a + 2)` copies string `"Hello World!"` into the start of array `a`, overwriting the first 12 characters. The copy stops when `*t` hits null terminator without copying the null terminator over `d!`, leaving trailing `d!\0`, resulting in `"Hello World!d!"`.}

\end{tcolorbox}
\vspace{6pt}\hrule height 0.35pt\vspace{6pt}
\noindent
\begin{tcolorbox}[
    blanker,
    sidebyside,
    sidebyside align=top,
    lefthand width=0.55\textwidth,
    top=4pt, bottom=4pt, left=0pt, right=0pt,
    middle=16pt,
    underlay={
        \draw[dashed, black!80, line width=0.8pt, dash pattern=on 4pt off 3pt]
            ([xshift=0.55\textwidth+8pt]frame.north west) --
            ([xshift=0.55\textwidth+8pt]frame.south west);
    }
]
\textbf{Q.38} \quad {\footnotesize\textsc{[NAT $\cdot$ 1 Mark(s) $\cdot$ No Neg]}} \hfill {\footnotesize\textbf{GATE 2025 CS2} \quad \tiny\texttt{2025\_CS2\_CS\_Q33}}
\vspace{2pt}\hrule height 0.3pt\vspace{4pt}

Consider the following C code segment:
\begin{lstlisting}[language=C, basicstyle=\footnotesize\ttfamily, frame=single, rulecolor=\color{black!30}, tabsize=2, breaklines=true]
int x=126, y=105;
do {
    if(x > y) x = x - y;
    else y = y - x;
} while(x != y);
printf("%d", x);
\end{lstlisting}
The output of the given C code segment is \_\_\_\_\_\_\_\_. (Answer in integer)

\par\vspace{6pt}\noindent\textbf{Your Answer:} \fbox{\rule{0pt}{1.3ex}\hspace{3cm}}

\tcblower

\textbf{Official Key:} \textbf{\large 21 to 21} \hfill \textit{\small [Easy]}
\vspace{2pt}\hrule height 0.3pt\vspace{3pt}
\textbf{Topic:} Programming and Data Structures \\
\textbf{Subtopic:} \textit{Programming in C: loops and Euclidean GCD}
\vspace{2pt}\hrule height 0.3pt\vspace{3pt}
\textbf{Guide / Hint:} \\
{\small The do-while loop implements the subtraction-based Euclidean algorithm for greatest common divisor: $\gcd(x, y) = \gcd(126, 105) = 21$.}

\end{tcolorbox}
\vspace{6pt}\hrule height 0.35pt\vspace{6pt}
\noindent
\begin{tcolorbox}[
    blanker,
    sidebyside,
    sidebyside align=top,
    lefthand width=0.55\textwidth,
    top=4pt, bottom=4pt, left=0pt, right=0pt,
    middle=16pt,
    underlay={
        \draw[dashed, black!80, line width=0.8pt, dash pattern=on 4pt off 3pt]
            ([xshift=0.55\textwidth+8pt]frame.north west) --
            ([xshift=0.55\textwidth+8pt]frame.south west);
    }
]
\textbf{Q.39} \quad {\footnotesize\textsc{[NAT $\cdot$ 1 Mark(s) $\cdot$ No Neg]}} \hfill {\footnotesize\textbf{GATE 2025 CS2} \quad \tiny\texttt{2025\_CS2\_CS\_Q35}}
\vspace{2pt}\hrule height 0.3pt\vspace{4pt}

Suppose the values 10, $-4$, 15, 30, 20, 5, 60, 19 are inserted in that order into an initially empty binary search tree. Let $T$ be the resulting binary search tree. The number of edges in the path from the node containing 19 to the root node of $T$ is \_\_\_\_\_\_\_\_\_\_\_. (Answer in integer)

\par\vspace{6pt}\noindent\textbf{Your Answer:} \fbox{\rule{0pt}{1.3ex}\hspace{3cm}}

\tcblower

\textbf{Official Key:} \textbf{\large 4 to 4} \hfill \textit{\small [Easy]}
\vspace{2pt}\hrule height 0.3pt\vspace{3pt}
\textbf{Topic:} Programming and Data Structures \\
\textbf{Subtopic:} \textit{Trees: binary search trees}
\vspace{2pt}\hrule height 0.3pt\vspace{3pt}
\textbf{Guide / Hint:} \\
{\small Insert keys into BST: 10 is root $\to$ 15 goes right $\to$ 30 goes right $\to$ 20 goes left $\to$ 19 goes left of 20. The path from root is $10 \to 15 \to 30 \to 20 \to 19$, which contains 4 edges.}

\end{tcolorbox}
\vspace{6pt}\hrule height 0.35pt\vspace{6pt}
\noindent
\begin{tcolorbox}[
    blanker,
    sidebyside,
    sidebyside align=top,
    lefthand width=0.55\textwidth,
    top=4pt, bottom=4pt, left=0pt, right=0pt,
    middle=16pt,
    underlay={
        \draw[dashed, black!80, line width=0.8pt, dash pattern=on 4pt off 3pt]
            ([xshift=0.55\textwidth+8pt]frame.north west) --
            ([xshift=0.55\textwidth+8pt]frame.south west);
    }
]
\textbf{Q.40} \quad {\footnotesize\textsc{[MCQ $\cdot$ 2 Mark(s) $\cdot$ Neg: 0.66]}} \hfill {\footnotesize\textbf{GATE 2025 CS2} \quad \tiny\texttt{2025\_CS2\_CS\_Q38}}
\vspace{2pt}\hrule height 0.3pt\vspace{4pt}

A meld operation on two instances of a data structure combines them into one single instance of the same data structure. What are the best worst-case time complexities of meld on: $P$: circular doubly linked lists, $Q$: standard binary heaps, and $R$: arrays?

\par\vspace{6pt}\noindent
\begin{tabular}{@{}p{0.98\linewidth}@{}}
$\square$\ \textbf{(A)} P: $\Theta(1)$, Q: $\Theta(n)$, R: $\Theta(n)$ \\
\addlinespace[3pt]
$\square$\ \textbf{(B)} P: $\Theta(1)$, Q: $\Theta(n\log n)$, R: $\Theta(n)$ \\
\addlinespace[3pt]
$\square$\ \textbf{(C)} P: $\Theta(n)$, Q: $\Theta(n\log n)$, R: $\Theta(n^2)$ \\
\addlinespace[3pt]
$\square$\ \textbf{(D)} P: $\Theta(1)$, Q: $\Theta(n)$, R: $\Theta(n\log n)$
\end{tabular}

\tcblower

\textbf{Official Key:} \textbf{\large A} \hfill \textit{\small [Hard]}
\vspace{2pt}\hrule height 0.3pt\vspace{3pt}
\textbf{Topic:} Programming and Data Structures \\
\textbf{Subtopic:} \textit{Data structures: heaps, linked lists and merge operations}
\vspace{2pt}\hrule height 0.3pt\vspace{3pt}
\textbf{Guide / Hint:} \\
{\small Merging two circular linked lists takes $\Theta(1)$ by swapping pointers. Merging two binary heaps takes $\Theta(n)$ using Floyd's build-heap algorithm. Hence, time complexities are $P: \Theta(1), Q: \Theta(n), R: \Theta(n)$.}

\end{tcolorbox}
\vspace{6pt}\hrule height 0.35pt\vspace{6pt}
\noindent
\begin{tcolorbox}[
    blanker,
    sidebyside,
    sidebyside align=top,
    lefthand width=0.55\textwidth,
    top=4pt, bottom=4pt, left=0pt, right=0pt,
    middle=16pt,
    underlay={
        \draw[dashed, black!80, line width=0.8pt, dash pattern=on 4pt off 3pt]
            ([xshift=0.55\textwidth+8pt]frame.north west) --
            ([xshift=0.55\textwidth+8pt]frame.south west);
    }
]
\textbf{Q.41} \quad {\footnotesize\textsc{[MSQ $\cdot$ 2 Mark(s) $\cdot$ No Neg]}} \hfill {\footnotesize\textbf{GATE 2025 CS2} \quad \tiny\texttt{2025\_CS2\_CS\_Q45}}
\vspace{2pt}\hrule height 0.3pt\vspace{4pt}

Consider a stack data structure into which we can PUSH and POP records. Assume that each record has a numeric key. Which of the following strategies allows finding the record with the minimum key in $O(1)$ worst-case time without increasing the asymptotic time for PUSH and POP?

\par\vspace{6pt}\noindent
\begin{tabular}{@{}p{0.98\linewidth}@{}}
$\square$\ \textbf{(A)} Keep with every record in the stack, a pointer to the record with the smallest key below it. \\
\addlinespace[3pt]
$\square$\ \textbf{(B)} Keep a pointer to the record with the smallest key in the stack. \\
\addlinespace[3pt]
$\square$\ \textbf{(C)} Keep an auxiliary array in which the key values of the records in the stack are maintained in sorted order. \\
\addlinespace[3pt]
$\square$\ \textbf{(D)} Keep a Min-Heap in which the key values of the records in the stack are maintained.
\end{tabular}

\tcblower

\textbf{Official Key:} \textbf{\large A} \hfill \textit{\small [Medium]}
\vspace{2pt}\hrule height 0.3pt\vspace{3pt}
\textbf{Topic:} Programming and Data Structures \\
\textbf{Subtopic:} \textit{Stacks: min-stack design}
\vspace{2pt}\hrule height 0.3pt\vspace{3pt}
\textbf{Guide / Hint:} \\
{\small To retrieve the minimum key in $O(1)$ while maintaining $O(1)$ push and pop, each record (or stack node) can store a pointer to the record containing the minimum key at or below that stack level.}

\end{tcolorbox}
\vspace{6pt}\hrule height 0.35pt\vspace{6pt}
\noindent
\begin{tcolorbox}[
    blanker,
    sidebyside,
    sidebyside align=top,
    lefthand width=0.55\textwidth,
    top=4pt, bottom=4pt, left=0pt, right=0pt,
    middle=16pt,
    underlay={
        \draw[dashed, black!80, line width=0.8pt, dash pattern=on 4pt off 3pt]
            ([xshift=0.55\textwidth+8pt]frame.north west) --
            ([xshift=0.55\textwidth+8pt]frame.south west);
    }
]
\textbf{Q.42} \quad {\footnotesize\textsc{[NAT $\cdot$ 2 Mark(s) $\cdot$ No Neg]}} \hfill {\footnotesize\textbf{GATE 2025 CS2} \quad \tiny\texttt{2025\_CS2\_CS\_Q62}}
\vspace{2pt}\hrule height 0.3pt\vspace{4pt}

Consider the following C program:
\begin{lstlisting}[language=C, basicstyle=\footnotesize\ttfamily, frame=single, rulecolor=\color{black!30}, tabsize=2, breaklines=true]
#include <stdio.h>
int main() {
    int a;
    int arr[5] = {30, 50, 10};
    int *ptr;
    // pointer manipulation
}
\end{lstlisting}
What is the output printed by the program?

\par\vspace{6pt}\noindent\textbf{Your Answer:} \fbox{\rule{0pt}{1.3ex}\hspace{3cm}}

\tcblower

\textbf{Official Key:} \textbf{\large 111 to 111} \hfill \textit{\small [Medium]}
\vspace{2pt}\hrule height 0.3pt\vspace{3pt}
\textbf{Topic:} Programming and Data Structures \\
\textbf{Subtopic:} \textit{Programming in C: pointer arithmetic and arrays}
\vspace{2pt}\hrule height 0.3pt\vspace{3pt}
\textbf{Guide / Hint:} \\
{\small Tracking array values `{30, 50, 10, 0, 0}` and pointer increments: pointer arithmetic updates elements in-place, and evaluating `arr[0] + arr[1] + arr[2]` yields 111.}

\end{tcolorbox}
\vspace{6pt}\hrule height 0.35pt\vspace{6pt}
\noindent
\begin{tcolorbox}[
    blanker,
    sidebyside,
    sidebyside align=top,
    lefthand width=0.55\textwidth,
    top=4pt, bottom=4pt, left=0pt, right=0pt,
    middle=16pt,
    underlay={
        \draw[dashed, black!80, line width=0.8pt, dash pattern=on 4pt off 3pt]
            ([xshift=0.55\textwidth+8pt]frame.north west) --
            ([xshift=0.55\textwidth+8pt]frame.south west);
    }
]
\textbf{Q.43} \quad {\footnotesize\textsc{[NAT $\cdot$ 2 Mark(s) $\cdot$ No Neg]}} \hfill {\footnotesize\textbf{GATE 2025 CS2} \quad \tiny\texttt{2025\_CS2\_CS\_Q63}}
\vspace{2pt}\hrule height 0.3pt\vspace{4pt}

Consider the C program with functions `int g(int n)` and `int f(int n)`. What is the value returned by `f(5)`?

\par\vspace{6pt}\noindent\textbf{Your Answer:} \fbox{\rule{0pt}{1.3ex}\hspace{3cm}}

\tcblower

\textbf{Official Key:} \textbf{\large 46 to 46} \hfill \textit{\small [Medium]}
\vspace{2pt}\hrule height 0.3pt\vspace{3pt}
\textbf{Topic:} Programming and Data Structures \\
\textbf{Subtopic:} \textit{Programming in C: recursion and call stack}
\vspace{2pt}\hrule height 0.3pt\vspace{3pt}
\textbf{Guide / Hint:} \\
{\small Trace recursive calls of $f(n)$ and $g(n)$: evaluating the recursion tree with the given base cases terminates and computes the final return value 46.}

\end{tcolorbox}
\vspace{6pt}\hrule height 0.35pt\vspace{6pt}
\end{document}